%% 物理第2巻『熱学』：組 熱W1（第22回編集委員会の手引き tools/BRIEF-phys2.md）
%% p2-01.tex  熱-1　熱の出入りを数える：温度・比熱・潜熱（★1〜4。3段にしない）
%% 記号（この章）：絶対温度 T〔K〕、セルシウス温度 θ〔℃〕（第22回の決定3。JIS Z 8202 の書き方。時刻 t とぶつけない）、
%%   温度差 ΔT〔K〕（℃ の差と同じ数）、熱量 Q〔J〕、電力・仕事率 P〔W〕、時刻 t〔s〕、熱容量 C_th〔J/K〕（数学 6-3・6-4 と同じ）、
%%   比熱 c〔J/(kg·K)〕、潜熱 L〔J/kg〕（融解熱 L_f・蒸発熱 L_v）、線膨張係数 α・体膨張率 β〔1/K〕、空気比 λ、
%%   体積流量 V-dot〔m^3/s〕（Q は熱量だけ）、質量流量 m-dot、ボルツマン定数 k_B（いつも添字）、効率 η。
%%   式の中の T はいつも K。℃ の値を使うときは θ と書く。
%% 数値はすべて Python で計算した：checks/物理第2巻/熱-1/check_p2_01.py（出力 check_p2_01_out.txt）
%%   水の比熱は本文 4.2 kJ/(kg·K)。3けたが要る例題1だけ 4.18（はじめに書いた）。
%% 第1巻・数学の本との約束：
%%   力-8 ⑦「線膨張係数（鋼 1.2・アルミ 2.3・銅 1.7 ×10^-5 /K）は熱-1 でもう一度使います」→ ⑦（同じ3つの値で β≈3α）・例題8。
%%   数学 6-3③（1 L の水 C_th=4200 J/K・1250 W）・6-4⑥（250 mL のお湯 1050 J/K・熱抵抗・時定数）・6-14⑧（80 ℃ 200 g＝840 J/K と 20 ℃ 600 g → 35 ℃）
%%   を ②③④・例題12 で同じ数のまま指す。数学第6巻が温度を T〔℃〕と書いたことと、この巻の T〔K〕・θ〔℃〕の関係は ② で説明した。
%% 出典（本文で使った事実・値）：
%%   ・2019年5月20日からの SI：ボルツマン定数 k_B＝1.380649×10^-23 J/K（定義の値）で K を決める。それまでは水の三重点 273.16 K で決めていた（BIPM SI 第9版）。
%%     セルシウス温度 θ/℃＝T/K−273.15（定義）。
%%   ・セルシウス（1742年に提案。はじめは水が凍る温度を100、わく温度を0とした）、ファーレンハイト（1724年）、ケルビン（W. トムソン、1848年に絶対温度の目もりを提案）。
%%   ・ジョゼフ・ブラック（1728〜1799年）：1760年ごろ、潜熱と比熱（熱容量）の考えを講義で示した。
%%   ・1 cal＝4.184 J（熱化学カロリー。計量法では、カロリーは人・動物が食べる物の熱量・代謝の熱量に限って使える）〔要確認：計量単位令の条文。二次資料で確認〕。
%%   ・水の比熱（20 ℃）4.184 kJ/(kg·K)、密度の最大 3.98 ℃（999.97 kg/m^3）：CoolProp 8.0（IAPWS-95）で計算。
%%   ・水の蒸発熱：0.01 ℃ 2501、25 ℃ 2442、35 ℃ 2418、100 ℃ 2257 kJ/kg：IAPWS-IF97 の式から計算（iapws 1.5.5）。
%%   ・氷の融解熱 334 kJ/kg、氷の比熱 2.1 kJ/(kg·K)、氷の密度 917 kg/m^3（0 ℃）：理科年表・Engineering ToolBox の値〔要確認：理科年表の最新版の本文〕。
%%   ・比熱：アルミニウム 0.90（897 J/(kg·K)、数学 8-13 と同じ）、鉄 0.45、銅 0.39（385）、エタノール 2.4、空気（圧力一定）1.0：CRC Handbook の値〔要確認〕。
%%     人の体の平均の比熱 約3.5 kJ/(kg·K)（生理学の教科書の目安。3.47 とする本が多い）〔要確認〕。
%%   ・ガソリンの体膨張率 約9.5×10^-4 /K（20 ℃前後。種類で変わる）〔要確認：JIS K 2249 の体積換算表〕。水 2.1×10^-4 /K（20 ℃、CoolProp で計算）。
%%   ・メタンの高発熱量 890.6 kJ/mol・低発熱量 802.6 kJ/mol（25 ℃）。水の蒸発熱 25 ℃で 44.0 kJ/mol（NIST の生成エンタルピー −285.83・−241.83 kJ/mol の差）。
%%     水素の高発熱量 285.8 kJ/mol・低発熱量 241.8 kJ/mol（同じ値から）。〔要確認：理科年表・JIS K 2279〕
%%   ・都市ガス 13A の標準熱量 45 MJ/m^3N（高発熱量。ガス事業者の供給約款の値）。
%%   ・ガス給湯器の「号数」：水温を 25 K 上げた湯を 1 分に何 L 出せるか（1号 ≒ 1.74 kW）。日本ガス石油機器工業会の説明〔要確認〕。
%%   ・氷の表面が滑りやすいわけ：R. Rosenberg, "Why is ice slippery?", Physics Today 58(12), 50–55 (2005)。
%%     圧力でとける説では −35 ℃でも滑ることを説明できない。摩擦の熱と、表面の液体のような薄い層の2つが考えられ、どちらがどれだけ効くかは研究が続いている。
%%   ・空気の酸素の割合 21 %（体積）。m^3N は 0 ℃・101.325 kPa の体積。
\makeatletter
% 温度計（棒）  \qxbthermo{(下の玉の中心)}{高さ}{液の高さ}{色}
\newcommand\qxbthermo[4]{\begin{scope}[shift={#1}]
  \filldraw[fill=white, draw=bkInk, line width=0.6pt, rounded corners=0.8mm] (-1.3,1.2) rectangle (1.3,#2);
  \fill[#4] (-0.6,0) rectangle (0.6,#3);
  \filldraw[fill=#4, draw=bkInk, line width=0.6pt] (0,0) circle (2.0);
\end{scope}}
% コップ  \qxbcup{(左下)}{はば}{高さ}{水の高さ}
\newcommand\qxbcup[4]{\begin{scope}[shift={#1}]
  \fill[bkBlue!18] (0.6,0.6) rectangle (#2-0.6,#4);
  \draw[bkInk, line width=0.8pt] (0,#3) -- (0,0) -- (#2,0) -- (#2,#3);
\end{scope}}
% 氷のかけら  \qxbice{(中心)}{大きさ}{回す角}
\newcommand\qxbice[3]{\filldraw[fill=white, draw=bkBlue!70!black, line width=0.6pt, rotate around={#3:#1}] ($#1+(-#2,-#2)$) rectangle ($#1+(#2,#2)$);}
% 燃焼の体積の単位 m^3_N（本文でも式でも使える）
\newcommand\qxbmN{\ensuremath{\,\mathrm{m}^3_{\mathrm N}}}
\makeatother
\tikzset{qxb lab/.style={font=\sffamily\footnotesize, inner sep=1pt},
  qxb arr/.style={-{Stealth[length=4pt, width=3.6pt]}, line width=1.4pt}}

\phchapter{1}{熱の出入りを数える：\\温度・比熱・潜熱}{氷水は、氷がとけきるまで0\,℃のままなのはなぜ？}{\bklevelrange{1}{4}}{熱の出入りを数える：温度・比熱・潜熱}
\bkmichishirube{2-3,2-10,2-16,6-3,6-4,6-14,力-8}{熱-2,熱-3,熱-8,熱-9,熱-11}

\begin{goal}
\begin{itemize}[label={\color{bkTheme}●}, leftmargin=1.6\zw, itemsep=0pt, topsep=1pt]
\item \bklevel{1}\enspace 氷がとけている間、温度が変わらないわけを、潜熱で説明できる。
\item \bklevel{2}\enspace $Q=mc\,\Delta T$ と熱のバランス式で、温まり方と混ぜたときの温度を求められる。
\item \bklevel{3}\enspace 加熱曲線を読み、熱膨張の大きさを見積もれる。
\item \bklevel{4}\enspace 燃料が燃えるときの空気の量と発熱量を、原子のバランスで求められる。
\end{itemize}
\end{goal}

\phkoushiki{%
\phk{物は温めると体積が増える（空気＞水＞金属）}{小学校理科}{ph:熱-1:bouchou}
\phk{水は約100\,℃でわき、0\,℃で氷になる。その間は温度が変わらない}{小学校理科}{ph:熱-1:santai}
\phk{熱は温度の高い物から低い物へ移り、同じ温度になる}{中学理科}{ph:熱-1:heikou}
\phk{状態変化：質量は変わらず、体積が変わる}{中学理科}{ph:熱-1:joutai}
\phk{熱量＝電力×時間 $Q=Pt$}{中学理科}{ph:熱-1:Pt}
\phk{熱運動と絶対温度 $T=\theta+273$}{高校物理基礎}{ph:熱-1:zettai}
\phk{熱容量 $C_{\mathrm{th}}=Q/\Delta T$}{高校物理基礎}{ph:熱-1:youryou}
\phk{比熱 $Q=mc\,\Delta T$}{高校物理基礎}{ph:熱-1:mcdT}
\phk{熱量の保存（混ぜたときの温度）}{高校物理基礎}{ph:熱-1:hozon}
\phk{融解熱・蒸発熱 $Q=mL$}{高校物理基礎}{ph:熱-1:mL}
\phk{体膨張率 $\beta\approx3\alpha$}{高校物理基礎}{ph:熱-1:3alpha}
}

\begin{hatena}
\noindent\begin{minipage}[c]{0.6\linewidth}\raggedright
コップに氷と水を入れ、温度計を入れておきます。部屋は25\,℃です。まわりから熱が入ってきて、氷は少しずつとけていきます。

ところが、温度計は、氷が残っているかぎり、ずっと0\,℃のままです。熱が入っているのに、温度が上がらないのはなぜでしょう。入った熱は、どこへ行ったのでしょう。
\end{minipage}\hfill
\begin{minipage}[c]{0.36\linewidth}\centering
\begin{tikzpicture}[x=1mm, y=1mm]
  \qxbcup{(0,0)}{30}{34}{26}
  \qxbice{(8,22)}{3.6}{12}\qxbice{(17,23)}{3.4}{-20}\qxbice{(24,21)}{3.2}{30}\qxbice{(12,15)}{3}{5}
  \qxbthermo{(20,5)}{42}{12}{bkRed}
  \node[qxb lab, anchor=west, text=bkRed!75!black] at (23,40) {0\,℃のまま};
  \node[qxb lab, anchor=north, text=bkGray] at (15,-1.5) {部屋は25\,℃};
\end{tikzpicture}
\end{minipage}
\end{hatena}

\noindent この巻は、熱の本です。熱で動くエンジン、熱を運ぶエアコン、熱を逃がさない家の壁。どれも、「熱がどれだけ入って、どれだけ出たか」を数えることから始まります。この章では、まず\textbf{温度}をきちんと決め、次に、物を温めるのに要る熱（比熱）と、とかしたり蒸発させたりするのに要る熱（潜熱）を数えます。数え方は、第1巻の力-13 と同じ「\textbf{入る量－出る量＝たまる量}」です。

\begin{yattemiyou}[氷水を温めると、温度はどう変わる？]
\textbf{用意するもの}\par
コップ、氷（コップの半分くらい）、水、温度計、ぬるま湯（40\,℃くらい）を入れたボウル、時計、記録の紙\par\vspace{2pt}
\noindent\begin{minipage}[c]{0.6\linewidth}
\begin{enumerate}[label={\bkcircled[bkGreen]{\arabic*}}, leftmargin=2.4\zw, itemsep=1pt, topsep=0pt]
\item コップに氷と水を入れ、よくかき混ぜてから温度を測る。
\item コップをぬるま湯のボウルにつけ、かき混ぜながら、30秒ごとに温度を記録する。
\item 氷がなくなってからも、5分くらい続ける。かき混ぜるのをやめると、どうなるかも見る。
\item 横に時間、たてに温度をとって、グラフにかく。
\end{enumerate}
\end{minipage}\hfill
\begin{minipage}[c]{0.36\linewidth}\centering
\begin{tikzpicture}[x=1mm, y=1mm]
  \fill[bkOrange!15] (0,0) rectangle (44,18);
  \draw[bkInk, line width=0.8pt] (0,24) -- (0,0) -- (44,0) -- (44,24);
  \qxbcup{(11,2)}{22}{30}{22}
  \qxbice{(17,18)}{2.8}{15}\qxbice{(25,19)}{2.6}{-15}
  \qxbthermo{(21,6)}{36}{9}{bkRed}
  \node[qxb lab, anchor=north] at (22,-1.5) {ぬるま湯につける};
\end{tikzpicture}
\end{minipage}\par\vspace{2pt}
{\small やけどをしないように、熱湯は使いません。温度計は、ガラスの棒ではなく、デジタルの温度計がわれにくく安全です。}
\end{yattemiyou}

\noindent やってみようの答え：氷が残っている間、温度はほぼ0\,℃のまま変わりません。氷がなくなると、温度は上がり始め、ぬるま湯の温度に近づいていきます。かき混ぜるのをやめると、温度計の所だけ少し温かくなることがあります。コップの中の温度が、場所で違ってしまうからです。この章では、「よくかき混ぜて、どこも同じ温度」の場合を考えます（どこも同じでないときの熱の伝わり方は、熱-11・熱-12 で学びます）。

\begin{omoidasou}
\bkrecall{2-3}{わり算で「1あたりの量」。$\text{熱容量}\div\text{質量}$ は「1\,kg あたりの熱容量」。}
\bkrecall{2-10}{単位の換算は「1をかける」。\qty{1}{kJ}＝\qty{1000}{J}、水\qty{1}{L}は約\qty{1}{kg}。}
\bkrecall{2-16}{1次方程式。混ぜたときの温度は、わからない温度を $x$ とおいて解く。}
\bkrecall{5-12}{電力を時間で積分すると、電力量（エネルギー）。1\,W は1秒あたり1\,J。}
\bkrecall{6-3}{熱容量 $C_{\mathrm{th}}$〔J/K〕。\qty{1}{L}の水は $C_{\mathrm{th}}=\qty{4200}{J/K}$。\qty{1250}{W}のヒーターで温めると $\dd T/\dd t=P/C_{\mathrm{th}}$。}
\bkrecall{6-4}{熱抵抗 $R_{\mathrm{th}}$〔K/W〕。逃げる熱は $\text{温度差}\div R_{\mathrm{th}}$、時定数は $\tau=C_{\mathrm{th}}R_{\mathrm{th}}$。}
\bkrecall{6-14}{2つの物の温度がそろうと、最後の温度は $\dfrac{C_1T_1+C_2T_2}{C_1+C_2}$。}
\bkrecall{力-8}{線膨張 $\Delta L=\alpha L\,\Delta T$。$\alpha$ は鋼 $1.2\times10^{-5}$、アルミニウム $2.3\times10^{-5}$、銅 $1.7\times10^{-5}$（1/K）。}
\end{omoidasou}

%% ======================== 物理の道具 ========================
\bksection{doguu}{物理の道具}

\bksubsection{1}{温度と熱平衡}
\label{ph:熱-1:sec1}
\phboxhead{熱いものと冷たいものをくっつけると}
熱いお茶の入った湯のみを、冷たい水の入ったおけにつけると、お茶は冷え、水は温まります。しばらくすると、どちらも同じ温度になって、それ以上は変わらなくなります。このとき、温度の高いほうから低いほうへ移ったものを\textbf{熱}といいます。\bkidx{ねつ}{熱}そして、くっつけた物どうしの温度がそろって、もう熱が移らない状態を\textbf{熱平衡}（ねつへいこう）といいます。\bkidx{ねつへいこう}{熱平衡}

\begin{bkfigure}
\begin{minipage}[c]{0.44\linewidth}\centering
\begin{tikzpicture}[x=1mm, y=1mm]
  \filldraw[fill=bkRed!25, draw=bkRed!80!black, line width=0.7pt, rounded corners=1mm] (0,0) rectangle (20,16);
  \filldraw[fill=bkBlue!20, draw=bkBlue!80!black, line width=0.7pt, rounded corners=1mm] (20,0) rectangle (40,16);
  \node[qxb lab] at (10,11) {高い温度};
  \node[qxb lab] at (30,11) {低い温度};
  \draw[qxb arr, bkOrange!85!black] (13,5) -- (27,5);
  \node[qxb lab, anchor=north, text=bkOrange!80!black] at (20,-1.5) {熱が移る};
\end{tikzpicture}
\end{minipage}\hfill
\begin{minipage}[c]{0.54\linewidth}\centering
\begin{tikzpicture}
\begin{axis}[bkplot, width=62mm, height=40mm, xmin=0, xmax=11, ymin=0, ymax=92, xtick={2,4,6,8,10}, ytick={20,35,80}, xlabel={時間（分）}, ylabel={温度（℃）}]
  \addplot[bkthin, dashed, domain=0:10.5] {35};
  \addplot[bkcurve2, domain=0:10.5, samples=60] {35+45*exp(-x/2)};
  \addplot[bkcurve, domain=0:10.5, samples=60] {35-15*exp(-x/2)};
  \node[qxb lab, anchor=south west, text=bkOrange!80!black] at (axis cs:1.6,62) {お茶};
  \node[qxb lab, anchor=north west, text=bkBlue!80!black] at (axis cs:1.6,24) {おけの水};
\end{axis}
\end{tikzpicture}
\end{minipage}
\bkfigcaption{熱平衡。熱は温度の高いほうから低いほうへ移り、2つの温度がそろうと止まる（そろい方の式は 6-14⑧）。}\label{ph:熱-1:fig-heikou}
\end{bkfigure}

\label{ph:熱-1:heikou}中学理科で学ぶ「熱は温度の高い物から低い物へ移り、やがて同じ温度になる」は、このことです。熱が自然に低い温度から高い温度へ移ることはありません。この「向き」のきまりは、熱-5 で第二法則として、くわしく学びます。

\phboxhead{温度計は、なぜ温度を測れる？：第0法則}
温度計を水に入れて、しばらく待つと、目もりが止まります。温度計と水が熱平衡になったのです。次に、同じ温度計を別の水に入れて、同じ目もりで止まったとします。このとき、2つの水をくっつけても、熱は移りません。経験から、次のことが知られています。
\begin{bkcenter}
\fbox{\parbox{0.86\linewidth}{\centering A と C が熱平衡で、B と C も熱平衡なら、A と B も熱平衡である。}}
\end{bkcenter}
これを\textbf{熱力学の第0法則}といいます。\bkidx{ねつりきがくのだいぜろほうそく}{熱力学の第0法則}C を温度計と考えると、温度計の目もりが同じ物どうしは、くっつけても熱が移らないとわかります。温度計が「同じ温度」を確かめる道具として使えるのは、この法則のおかげです。理屈で導いたのではなく、どんな実験でも成り立ってきた\textbf{経験則}です。第1法則・第2法則よりあとに、その前に置くべき大切な法則だと気づかれたので、「第0」という番号がつきました。

\phboxhead{温度計のしくみ}
温度計は、「温度で変わる性質」を目もりにした道具です。
\begin{bkfigure}
\begin{tikzpicture}[x=1mm, y=1mm]
  % 液柱
  \qxbthermo{(10,4)}{38}{24}{bkRed}
  \foreach \yy in {10,16,22,28,34} {\draw[bkInk, line width=0.4pt] (11.3,\yy) -- (13.3,\yy);}
  \node[qxb lab, align=center, anchor=north] at (10,-1) {液柱温度計\\液がふくらむ（⑦）};
  % 熱電対
  \draw[bkOrange!85!black, line width=1.2pt] (40,4) -- (52,30) -- (58,30);
  \draw[bkBlue!80!black, line width=1.2pt] (40,4) -- (60,26) -- (66,26);
  \fill[bkInk] (40,4) circle (1.0);
  \filldraw[fill=bkGray!20, draw=bkInk, line width=0.6pt, rounded corners=1mm] (57,22) rectangle (75,36);
  \node[qxb lab] at (66,31.5) {mV};
  \node[qxb lab, align=center, anchor=north] at (54,-1) {熱電対\\接点の電圧（電-4）};
  % 放射温度計
  \filldraw[fill=bkGray!25, draw=bkInk, line width=0.6pt, rounded corners=1mm] (98,18) -- (116,22) -- (116,28) -- (98,30) -- cycle;
  \filldraw[fill=bkGray!40, draw=bkInk, line width=0.6pt] (108,18.5) rectangle (113,8);
  \filldraw[fill=bkRed!30, draw=bkRed!80!black, line width=0.6pt] (84,8) rectangle (90,38);
  \foreach \yy in {18,24,30} {\draw[bkRed, line width=0.8pt, dashed, -{Stealth[length=3pt]}] (91,\yy) -- (97.5,{24+(\yy-24)*0.4});}
  \node[qxb lab, align=center, anchor=north] at (101,-1) {放射温度計\\出てくる赤外線（熱-13）};
\end{tikzpicture}
\bkfigcaption{3つの温度計。どれも、温度で変わる性質（体積・電圧・光の強さ）を目もりにしている。}\label{ph:熱-1:fig-ondokei}
\end{bkfigure}
\textbf{液柱温度計}は、細い管の中の液（色をつけた灯油など）が温まってふくらむのを、長さで読みます（⑦）。\textbf{熱電対}（ねつでんつい）は、2種類の金属線の先をつないだもので、つないだ所の温度で小さな電圧（数 mV）が出ます。工場の炉やエンジンの温度を測るのに使います（しくみは物理シリーズの『電磁気学』（電-4））。\textbf{放射温度計}は、物から出る目に見えない光（赤外線）の強さで、さわらずに温度を測ります。体温を額で測る温度計も、このなかまです（熱-13）。

\begin{phmijika}[温度と熱は、違うもの]
おふろのお湯と、おふろからくんだコップ1ぱいのお湯は、同じ40\,℃です。でも、冷えた部屋を温める力は、おふろのほうがずっと大きいでしょう。\textbf{温度}は「熱い・冷たいの度合い」で、量によりません。\textbf{熱}は「移っていくエネルギーの量」で、同じ温度でも、量が多いほど多く出し入れできます。この2つを区別するのが、この章のはじめの一歩です。
\end{phmijika}

%% ---------------------------------------------------------------
\bksubsection{2}{温度の目もり：℃ と K}
\label{ph:熱-1:sec2}
\phboxhead{セルシウス温度 $\theta$〔℃〕}
ふだん使う温度の目もりは、1気圧で水が凍る温度を0\,℃、わく温度を100\,℃にして、その間を100に分けたものです。スウェーデンのセルシウスが1742年に提案したので、\textbf{セルシウス温度}（摂氏温度）といいます。\bkidx{せるしうすおんど}{セルシウス温度}この本では、セルシウス温度を、ギリシャ文字の $\theta$（シータ）で表します。\textbf{時刻 $t$ とぶつからないようにするため}です（JIS の量と単位の決まり JIS Z 8202 も、セルシウス温度に $\theta$ か $t$ を使います）。

\phboxhead{絶対温度 $T$〔K〕}
物を冷やしていくと、どこまでも冷えるわけではありません。$-273.15$\,℃ より低い温度はない、ということがわかっています（気体の実験から、この温度が見つかったしくみは熱-3 ②）。このいちばん低い温度を\textbf{絶対零度}といい、そこを0にして、セルシウス温度と同じ目もりの大きさで数えた温度を\textbf{絶対温度}（熱力学温度）といいます。\bkidx{ぜったいおんど}{絶対温度}単位は K（ケルビン）で、記号は $T$ です。\bkidx{けるびん}{ケルビン（K）}
\pheq{T=\theta+273.15}{K ＝ ℃の数＋273.15}
\label{ph:熱-1:zettai}高校の物理基礎では、273.15 を 273 にまるめて $T=\theta+273$ と書きます。たとえば20\,℃は約\qty{293}{K}、絶対零度は\qty{0}{K}、水がわく100\,℃は約\qty{373}{K}です。

\begin{phyomi}[$T=\theta+273.15$]
「目もりの間かくは℃と同じ。0の位置だけを、273.15 だけ下へずらした目もり」。だから、絶対温度は負の数になりません。
\end{phyomi}

\begin{bkfigure}
\begin{tikzpicture}[x=1mm, y=0.135mm]
  \foreach \xx/\nm in {10/K, 42/℃, 74/$^{\circ}$F} {\draw[bkInk, line width=1pt] (\xx,0) -- (\xx,400); \node[qxb lab, anchor=south, font=\sffamily\small] at (\xx,410) {\nm};}
  \foreach \tt/\lab in {0/絶対零度, 273.15/水が凍る, 310.15/体温, 373.15/水がわく} {
    \draw[bkGray, line width=0.4pt, dashed] (6,\tt) -- (86,\tt);
    \node[qxb lab, anchor=west, text=bkInk] at (88,\tt) {\lab};}
  \node[qxb lab, anchor=east, fill=bkPaper] at (9,0) {0};       \node[qxb lab, anchor=east, fill=bkPaper] at (41,0) {$-273.15$};  \node[qxb lab, anchor=east, fill=bkPaper] at (73,0) {$-459.67$};
  \node[qxb lab, anchor=east, fill=bkPaper] at (9,273.15) {273.15}; \node[qxb lab, anchor=east, fill=bkPaper] at (41,273.15) {0};   \node[qxb lab, anchor=east, fill=bkPaper] at (73,273.15) {32};
  \node[qxb lab, anchor=east, fill=bkPaper] at (9,310.15) {310.15}; \node[qxb lab, anchor=east, fill=bkPaper] at (41,310.15) {37};  \node[qxb lab, anchor=east, fill=bkPaper] at (73,310.15) {98.6};
  \node[qxb lab, anchor=east, fill=bkPaper] at (9,373.15) {373.15}; \node[qxb lab, anchor=east, fill=bkPaper] at (41,373.15) {100}; \node[qxb lab, anchor=east, fill=bkPaper] at (73,373.15) {212};
\end{tikzpicture}
\bkfigcaption{3つの温度の目もり（1気圧の水で）。K と℃は目もりの間かくが同じで、0の位置だけが違う。}\label{ph:熱-1:fig-memori}
\end{bkfigure}

\phboxhead{温度の差は、K でも℃でも同じ数}
ある物の温度が $\theta_1$ から $\theta_2$ に変わったとします。絶対温度で書いた差は
\[\Delta T=T_2-T_1=(\theta_2+273.15)-(\theta_1+273.15)=\theta_2-\theta_1\]
で、273.15 が消えます。\textbf{温度の差は、K で数えても℃で数えても、同じ数}です。20\,℃から35\,℃に上がったら、「\qty{15}{K}上がった」と言っても「15\,℃上がった」と言っても同じです。この本では、温度の差は $\Delta T$ と書き、単位は K を使います。\qty{1}{K}と1\,℃は同じ大きさです。

\phboxhead{数学第6巻の $T$〔℃〕との関係}
数学第6巻（6-3・6-4・6-14）では、温度を $T$〔℃〕と書いて、「お湯が冷める」「ヒーターで温める」を計算しました。その式は、どれも、
\[C_{\mathrm{th}}\frac{\dd T}{\dd t}=P,\qquad \frac{\dd T}{\dd t}=-\frac{T-T_{\mathrm a}}{C_{\mathrm{th}}R_{\mathrm{th}}},\qquad T_\infty=\frac{C_1T_1+C_2T_2}{C_1+C_2}\]
の形でした。1つめと2つめには、温度の\textbf{変わる速さ}と\textbf{差}しか出てきません。3つめは「重みつきの平均」です。全部の温度に273.15をたすと、分子は $(C_1+C_2)\times273.15$ だけ増え、それを分母 $C_1+C_2$ でわるので、答えもちょうど273.15だけ大きくなります。だから、これらの式は、℃で入れても K で入れても正しく、数学の本の書き方で問題はありませんでした。

ところが、熱学には、\textbf{温度そのものの値}が効く式がたくさんあります。気体の状態方程式 $pV=nRT$（熱-3）、熱機関の効率の限界 $1-T_{\mathrm L}/T_{\mathrm H}$（熱-5）、熱ふく射の $T^4$（熱-13）などです。これらは K で入れないと、間違った答えになります。そこで、この巻では次のように決めます。
\begin{bkcenter}
\fbox{\parbox{0.86\linewidth}{\raggedright
・式の中の $T$ は、いつも絶対温度〔K〕。\\
・℃で書いた値を使うときは $\theta$〔℃〕と書き、式に入れる前に $T=\theta+273.15$ に直す。\\
・温度の差 $\Delta T$ は、K と℃で同じ数なので、どちらで計算してもよい（単位は K と書く）。}}
\end{bkcenter}
「20\,℃から40\,℃にしたら、2倍熱い」とは言えません。絶対温度では\qty{293}{K}から\qty{313}{K}で、約1.07倍です。

\phboxhead{温度の正体：分子の熱運動}
水も空気も、目に見えない小さな分子が集まってできています。分子は、止まっているのではなく、いつもでたらめな向きに動いたり、ふるえたりしています。この動きを\textbf{熱運動}といいます。\bkidx{ねつうんどう}{熱運動}温度が高いほど、熱運動ははげしくなります。絶対零度は、熱運動がいちばん小さくなる温度です。

2019年5月から、国際単位系（SI）では、\textbf{ボルツマン定数} $k_{\mathrm B}$ を次の値に決めることで、ケルビンを決めています。\bkidx{ぼるつまんていすう}{ボルツマン定数}
\[k_{\mathrm B}=\qty{1.380649e-23}{J/K}\qquad\text{（決めた値。誤差はない）}\]
それまでは、「水の三重点（氷・水・水蒸気がいっしょにある温度）を\qty{273.16}{K}とする」と決めていました。$k_{\mathrm B}$ は、温度 \qty{1}{K} が分子1個あたりどれくらいのエネルギーにあたるかを表す数です。分子の速さと温度の関係は、熱-14 で導きます。

\begin{seikaku}
\phboxhead{華氏温度（$^{\circ}$F）}
アメリカでは、ファーレンハイトが1724年に提案した\textbf{華氏温度}を使います。水が凍る温度が32\,$^{\circ}$F、わく温度が212\,$^{\circ}$F で、その間を180に分けています。セルシウス温度 $\theta$ との関係は
\[\theta_{\mathrm F}=1.8\,\theta+32\qquad\text{（$\theta_{\mathrm F}$ は華氏の値）}\]
です。体温の98.6\,$^{\circ}$F は、$(98.6-32)\div1.8=37.0\,℃$ です。華氏の1目もりは $\tfrac{100}{180}=\tfrac{5}{9}$\,K で、K・℃の目もりより小さいので、温度の差は同じ数になりません。
\end{seikaku}

%% ---------------------------------------------------------------
\bksubsection{3}{熱量・熱容量・比熱}
\label{ph:熱-1:sec3}
\phboxhead{熱はエネルギー：単位は J}
熱は、温度の差で移っていく\textbf{エネルギー}です。だから、熱の量（\textbf{熱量}）は、仕事や運動エネルギーと同じ J（ジュール）で数えます。\bkidx{ねつりょう}{熱量}記号は $Q$ です。熱がエネルギーであることは、仕事で物が温まることから確かめられました（熱-2 ①）。

昔は、「水\qty{1}{g}の温度を1\,℃上げる熱」をもとにした cal（カロリー）という単位を使っていました。今は $\qty{1}{cal}=\qty{4.184}{J}$ と決めてあり、日本の計量法では、カロリーは食べ物の熱量（栄養）を表すときだけに使えます。

\phboxhead{電気で温める：$Q=Pt$}
電気ポットやヒーターに入れた電気のエネルギーは、ほとんど全部が熱になります。電力 $P$〔W〕は「1秒あたりのエネルギー」なので、時刻 $0$ から $t$ までに入る熱は、電力を時間で積分した量（5-12⑥）です。
\pheq{Q=\int_0^t P\,\dd t}{J ＝ W × s}
電力が一定なら、長方形の面積で
\pheq{Q=Pt}{J ＝ W × s}
\label{ph:熱-1:Pt}です。中学理科で学ぶ「熱量〔J〕＝電力〔W〕×時間〔s〕」です。\qty{1200}{W}のヒーターを\qty{60}{s}使うと、$1200\times60=\qty{72000}{J}=\qty{72}{kJ}$ の熱が出ます。電流が熱を出すしくみ（ジュール熱）は、物理シリーズの『電磁気学』（電-4）で学びます。

\phboxhead{熱容量：1\,K 上げるのに要る熱}
同じ熱を入れても、温度の上がり方は物で違います。物の温度を\qty{1}{K}上げるのに要る熱を\textbf{熱容量}といい、数学 6-3③ と同じ $C_{\mathrm{th}}$ で書きます（th は「熱の」thermal の頭文字）。\bkidx{ねつようりょう}{熱容量}熱 $Q$ を入れて温度が $\Delta T$ 上がったなら、
\pheq{C_{\mathrm{th}}=\frac{Q}{\Delta T}}{J/K ＝ J ÷ K}
\label{ph:熱-1:youryou}です。6-3③ では、\qty{1}{L}の水が $C_{\mathrm{th}}=\qty{4200}{J/K}$ で、\qty{1250}{W}のヒーターなら1秒に $1250\div4200\approx\qty{0.298}{K}$ ずつ上がりました。熱容量が大きい物ほど、温まりにくく、冷めにくい物です。

\phboxhead{比熱：1\,kg あたりの熱容量}
同じ物でも、量が2倍なら熱容量も2倍です。物の種類だけで決まる数にするために、質量 $m$ でわった「1\,kg あたりの熱容量」を考えます（2-3③）。これを\textbf{比熱}（比熱容量）といい、$c$ で書きます。\bkidx{ひねつ}{比熱}
\pheq{c=\frac{C_{\mathrm{th}}}{m}}{J/(kg·K) ＝ (J/K) ÷ kg}
単位は J/(kg·K) です。水の比熱は約 4.2\,kJ/(kg·K) で、これは「水\qty{1}{kg}を\qty{1}{K}上げるのに約\qty{4.2}{kJ}」、言いかえると「水\qty{1}{g}を\qty{1}{K}上げるのに約\qty{4.2}{J}」（6-4⑥と同じ数）です。

\begin{bkcenter}
\begin{bktable}{colspec={Q[l,m,wd=40mm] Q[c,m,wd=38mm] Q[l,m,wd=42mm]}, rowsep=1.2pt, row{1}={bg=bkTheme!10}}
物 & 比熱 $c$〔kJ/(kg·K)〕 & 身近な例 \\
水（20\,℃） & 4.18（本文は 4.2） & 海・おふろ・冷却水 \\
人の体（平均） & 約 3.5 & 汗で体温を下げる \\
エタノール & 2.4 & 消毒用のアルコール \\
氷（0\,℃付近） & 2.1 & 保冷剤・かき氷 \\
空気（圧力一定） & 1.0 & 部屋の空気・エアコン \\
アルミニウム & 0.90 & なべ・放熱板 \\
鉄 & 0.45 & 鉄なべ・ブレーキの円板 \\
銅 & 0.39 & 電線・銅なべ \\
\end{bktable}
\bktabcaption{いろいろな物の比熱（およその値。巻末の物性値の表）。水の比熱は、とても大きい。}\label{ph:熱-1:tab-hinetsu}
\end{bkcenter}

\phboxhead{比熱から熱量を出す：$Q=mc\,\Delta T$}
熱容量は $C_{\mathrm{th}}=mc$ なので、温度を $\Delta T$ 上げるのに要る熱は $Q=C_{\mathrm{th}}\Delta T=mc\,\Delta T$ です。ただし、比熱は温度で少し変わります（図\ref{ph:熱-1:fig-cw}）。温度を少し $\dd T$ 上げるのに要る熱は $mc(T)\,\dd T$ なので、$T_1$ から $T_2$ まで温める熱は、それをたした積分です。
\pheq{Q=\int_{T_1}^{T_2}m\,c(T)\,\dd T}{J ＝ kg × J/(kg·K) × K}
比熱がほぼ一定とみなせる温度のはばなら、$c$ を積分の外に出して
\pheq{Q=mc\,\Delta T}{J ＝ kg × J/(kg·K) × K}
\label{ph:熱-1:mcdT}です。これが、高校の物理基礎の「熱量＝質量×比熱×温度の変化」です。

\begin{bkfigure}
\begin{tikzpicture}
\begin{axis}[bkplot, width=104mm, height=50mm, xmin=0, xmax=104, ymin=4.16, ymax=4.235, xtick={20,40,60,80,100}, ytick={4.18,4.2,4.22}, yticklabels={4.18,4.20,4.22}, xlabel={$\theta$（℃）}, ylabel={$c$（kJ/(kg·K)）}, axis lines=left]
  \addplot[bkcurve, smooth] coordinates {(0.1,4.2191) (5,4.2049) (10,4.1952) (15,4.1885) (20,4.1841) (25,4.1813) (30,4.1798) (35,4.1793) (40,4.1794) (50,4.1813) (60,4.1850) (70,4.1901) (80,4.1968) (90,4.2052) (99.9,4.2156)};
  \node[qxb lab, anchor=south, text=bkBlue!80!black, fill=bkPaper] at (axis cs:52,4.207) {0〜100\,℃で、ずれは1\,\% より小さい};
\end{axis}
\end{tikzpicture}
\bkfigcaption{1気圧の水の比熱（CoolProp 8.0 で計算）。ほぼ一定なので、$Q=mc\,\Delta T$ の $c$ を一定にしてよい。}\label{ph:熱-1:fig-cw}
\end{bkfigure}

\begin{phflow}[熱量]
\phflowrow{$Q=\displaystyle\int_0^t P\,\dd t$}{電力 $P$ が一定}{$Q=Pt$}
\phflowrow{$Q=\displaystyle\int_{T_1}^{T_2}m\,c(T)\,\dd T$}{比熱 $c$ が一定\\（状態変化なし）}{$Q=mc\,\Delta T$}
\end{phflow}

\begin{phyomi}[$Q=mc\,\Delta T$]
「温度を上げるのに要る熱は、量（$m$）と、物の温まりにくさ（$c$）と、上げる温度（$\Delta T$）のかけ算」。$\Delta T$ は差なので、K でも℃でも同じ数を入れてよい。
\end{phyomi}

\begin{phbaisuru}{水の質量 $m$}
同じ熱 $Q$ を入れたときの温度の上がり $\Delta T=Q/(mc)$ は半分になります。同じヒーターでお湯をわかすと、水の量が2倍なら、かかる時間も2倍です。
\par\smallskip\noindent\begin{tikzpicture}
\begin{axis}[bkplot, width=96mm, height=42mm, xmin=0, xmax=460, ymin=0, ymax=105, xtick={100,200,300,400}, ytick={25,50,100}, xlabel={$Q$（kJ）}, ylabel={$\Delta T$（K）}]
  \addplot[bkcurve, domain=0:440] {x/4.2};
  \addplot[bkcurve2, domain=0:440] {x/8.4};
  \node[qxb lab, anchor=south east, text=bkBlue!80!black] at (axis cs:330,80) {水\qty{1}{kg}};
  \node[qxb lab, anchor=north west, text=bkOrange!80!black] at (axis cs:330,30) {水\qty{2}{kg}};
  \fill[bkRed] (axis cs:210,50) circle (1.8pt) (axis cs:210,25) circle (1.8pt);
  \node[qxb lab, anchor=east, text=bkRed!75!black, fill=bkPaper] at (axis cs:202,37) {同じ\qty{210}{kJ}で半分};
\end{axis}
\end{tikzpicture}
\end{phbaisuru}

\begin{phmijika}[海辺の気温は、なぜ変わりにくい？]
水の比熱は、岩や砂（約 0.8\,kJ/(kg·K)）の5倍ほどもあります。夏の昼、同じ日ざしを受けても、砂浜は熱くなるのに、海の水はあまり温まりません。夜は、砂浜はすぐ冷えますが、海はなかなか冷えません。だから、海の近くの町は、内陸の町より、1日の気温の変化が小さくなります。車のエンジンを水で冷やすのも、水が少しの温度の上がりでたくさんの熱を運べるからです。
\end{phmijika}

%% ---------------------------------------------------------------
\bksubsection{4}{熱のバランス式と混ぜたときの温度}
\label{ph:熱-1:sec4}
\phboxhead{熱のバランス式}
物（ひとまとまりの物を\textbf{系}といいます。熱-2 ②）に入る熱と出る熱を数えます。力-13 と同じく、
\pheq{\text{たまる熱}=\text{入る熱}-\text{出る熱}}{J ＝ J － J}
です。\bkidx{ねつのばらんすしき}{熱のバランス式}たまった熱の分だけ、温度が上がります（$C_{\mathrm{th}}\Delta T$）。6-4⑥ の式
\[C_{\mathrm{th}}\frac{\dd T}{\dd t}=P-\frac{T-T_{\mathrm a}}{R_{\mathrm{th}}}\]
は、このバランス式を1秒あたりで書いたものでした（入る熱が $P$、逃げる熱が右の項）。

\phboxhead{外と熱をやりとりしないとき}
いくつかの物を、外と熱をやりとりしない入れ物（断熱した入れ物）の中でくっつけます。外からは入らず、外へも出ないので、全体にたまる熱は0です。1つ1つの物の熱容量を $C_1,\ C_2,\ \dots$、温度の変化を $\Delta T_1,\ \Delta T_2,\ \dots$ とすると、
\pheq{C_1\,\Delta T_1+C_2\,\Delta T_2+\cdots=0}{J ＝ J}
です。温度が下がった物の $\Delta T$ は負の数なので、「温度が下がった物が出した熱＝温度が上がった物が受けとった熱」と言っても同じです。これを\textbf{熱量の保存}といいます。\bkidx{ねつりょうのほぞん}{熱量の保存}

\phboxhead{2つの物を混ぜる}
温度 $T_1$ の物（熱容量 $C_1$）と、温度 $T_2$ の物（熱容量 $C_2$）を混ぜて、最後に同じ温度 $T_{\mathrm f}$ になったとします。$\Delta T_1=T_{\mathrm f}-T_1$、$\Delta T_2=T_{\mathrm f}-T_2$ なので、
\[C_1(T_{\mathrm f}-T_1)+C_2(T_{\mathrm f}-T_2)=0\]
これは $T_{\mathrm f}$ の1次方程式（2-16）です。$T_{\mathrm f}$ をふくむ項を左に集めて、
\pheq{T_{\mathrm f}=\frac{C_1T_1+C_2T_2}{C_1+C_2}}{K ＝ (J/K×K)÷(J/K)}
\label{ph:熱-1:hozon}です。6-14⑧ で、2つの温度がそろっていく連立微分方程式から求めた最後の温度と、同じ式です（混ぜても、金属の壁ごしでも、熱の合計は変わらないから）。②で見たように、この式は重みつきの平均なので、℃で入れても答えは同じ形で出ます。
\[\theta_{\mathrm f}=\frac{C_1\theta_1+C_2\theta_2}{C_1+C_2}\]
同じ物（たとえば水と水）なら、$C=mc$ の $c$ が約分されて、
\[\theta_{\mathrm f}=\frac{m_1\theta_1+m_2\theta_2}{m_1+m_2}\]
です。6-14⑧ の例題（80\,℃のお湯\qty{200}{g}と20\,℃の水\qty{600}{g}）なら、
\[\theta_{\mathrm f}=\frac{200\times80+600\times20}{800}=35\ \text{℃}\]
です。

\begin{phflow}[混ぜたときの温度]
\phflowrow{$\textstyle\sum_i C_i\,\Delta T_i=0$\\（外と熱をやりとりしない）}{物が2つ\\最後に同じ温度}{$T_{\mathrm f}=\dfrac{C_1T_1+C_2T_2}{C_1+C_2}$}
\end{phflow}

\begin{phyomi}[$T_{\mathrm f}=\dfrac{C_1T_1+C_2T_2}{C_1+C_2}$]
「最後の温度は、熱容量で重みをつけた平均」。てんびんの支点の位置と同じで、熱容量の大きいほう（重いほう）の温度に近くなる。
\end{phyomi}

\begin{bkfigure}
\begin{tikzpicture}[x=1mm, y=1mm]
  \draw[bkInk, line width=1.2pt] (0,10) -- (100,10);
  \foreach \xx/\tt/\yy in {0/20/7.5, 25/35/3, 100/80/7.5} {\draw[bkInk, line width=0.5pt] (\xx,8.5) -- (\xx,11.5); \node[qxb lab, anchor=north] at (\xx,\yy) {\tt\,℃};}
  \fill[bkInk] (25,10) -- (22,4) -- (28,4) -- cycle;
  \filldraw[fill=bkBlue!25, draw=bkBlue!80!black, line width=0.7pt] (-9,11) rectangle (9,26);
  \node[qxb lab, align=center] at (0,18.5) {水\qty{600}{g}\\\qty{2520}{J/K}};
  \filldraw[fill=bkRed!25, draw=bkRed!80!black, line width=0.7pt] (94,11) rectangle (106,17);
  \node[qxb lab, anchor=south, align=center] at (100,17.5) {お湯\qty{200}{g}\\\qty{840}{J/K}};
  \draw[{Stealth[length=3pt]}-{Stealth[length=3pt]}, bkGray, line width=0.5pt] (0,30) -- node[qxb lab, above] {15} (25,30);
  \draw[{Stealth[length=3pt]}-{Stealth[length=3pt]}, bkGray, line width=0.5pt] (25,30) -- node[qxb lab, above] {45} (100,30);
\end{tikzpicture}
\bkfigcaption{混ぜた温度は、てんびんのつり合う所。$2520\times15=840\times45$（力-5 のモーメントのつり合いと同じ形）。}\label{ph:熱-1:fig-tenbin}
\end{bkfigure}

\begin{phbaisuru}{冷たい水の量}
80\,℃のお湯\qty{200}{g}に、20\,℃の水を混ぜます。水の量を増やすほど、最後の温度は20\,℃に近づきますが、20\,℃より下がることはありません。水が\qty{200}{g}なら50\,℃、2倍の\qty{400}{g}なら40\,℃、さらに2倍の\qty{800}{g}なら32\,℃です。
\par\smallskip\noindent\begin{tikzpicture}
\begin{axis}[bkplot, width=96mm, height=50mm, xmin=0, xmax=1250, ymin=0, ymax=88, xtick={200,400,800,1200}, ytick={20,40,50,80}, xlabel={水の質量（g）}, ylabel={$\theta_{\mathrm f}$（℃）}]
  \addplot[bkthin, dashed, domain=0:1200] {20};
  \addplot[bkcurve, domain=0:1200, samples=80] {(200*80+x*20)/(200+x)};
  \fill[bkRed] (axis cs:200,50) circle (1.8pt) (axis cs:400,40) circle (1.8pt) (axis cs:800,32) circle (1.8pt);
\end{axis}
\end{tikzpicture}
\end{phbaisuru}

\begin{phmodel}[混ぜたときの温度]
\phignore{入れ物や温度計が受けとる熱。まわりの空気へ逃げる熱。お湯から出る湯気（蒸発の熱。⑤）。}
\phassume{最後は全体が同じ温度になる。比熱は温度によらず一定。}
\phlimit{入れ物の熱容量が中身にくらべて無視できないとき（⑨）。混ぜるのに時間がかかり、その間に熱が逃げるとき。}
\end{phmodel}

%% ---------------------------------------------------------------
\bksubsection{5}{状態変化と潜熱}
\label{ph:熱-1:sec5}
\phboxhead{固体・液体・気体}
水は、冷やすと氷（固体）になり、温めると水蒸気（気体）になります。このように、物が固体・液体・気体の間で姿を変えることを\textbf{状態変化}といいます。\bkidx{じょうたいへんか}{状態変化}固体がとけて液体になる温度を\textbf{融点}、液体がわき立って気体になる温度を\textbf{沸点}といいます。1気圧の水では、融点が0\,℃、沸点が約100\,℃です（1気圧で 99.97\,℃）。

\begin{bkfigure}
\begin{tikzpicture}[x=1mm, y=1mm]
  % 固体
  \draw[bkInk, line width=0.7pt, rounded corners=1mm] (0,0) rectangle (32,26);
  \foreach \i in {0,...,4} {\foreach \j in {0,...,3} {\fill[bkBlue!70] (4+\i*6,4+\j*6) circle (1.6);}}
  \foreach \i in {0,...,3} {\foreach \j in {0,...,3} {\draw[bkBlue!50, line width=0.6pt] (4+\i*6,4+\j*6) -- (10+\i*6,4+\j*6);}}
  \node[qxb lab, anchor=north, align=center] at (16,-1.5) {固体（氷）\\きちんと並び、手をつなぐ};
  % 液体
  \draw[bkInk, line width=0.7pt, rounded corners=1mm] (44,0) rectangle (76,26);
  \foreach \px/\py in {48/4,53/5,59/4,65/5,71/4,50/9,56/10,62/9,68/10,47/14,53/15,60/14,66/15,72/14,51/19.5,57/20,63/19} {\fill[bkBlue!70] (\px,\py) circle (1.6);}
  \node[qxb lab, anchor=north, align=center] at (60,-1.5) {液体（水）\\ふれ合いながら動く};
  % 気体
  \draw[bkInk, line width=0.7pt, rounded corners=1mm] (88,0) rectangle (120,26);
  \foreach \px/\py/\ang in {93/6/30,104/18/200,113/9/120,98/21/300,116/22/240} {\fill[bkBlue!70] (\px,\py) circle (1.6); \draw[bkGray, -{Stealth[length=2.4pt]}, line width=0.6pt] (\px,\py) -- ++(\ang:5);}
  \node[qxb lab, anchor=north, align=center] at (104,-1.5) {気体（水蒸気）\\ばらばらに飛ぶ};
  \draw[qxb arr, bkRed!80!black] (33,16) -- node[qxb lab, above] {融解} (43,16);
  \draw[qxb arr, bkRed!80!black] (77,16) -- node[qxb lab, above] {蒸発} (87,16);
  \draw[qxb arr, bkBlue!80!black] (43,9) -- node[qxb lab, below] {凝固} (33,9);
  \draw[qxb arr, bkBlue!80!black] (87,9) -- node[qxb lab, below] {凝縮} (77,9);
\end{tikzpicture}
\bkfigcaption{状態変化と分子の並び方（分子は大きくかいた）。右へ進む向き（矢印が右向き）は熱を入れる変化、左へ進む向きは熱を出す変化。}\label{ph:熱-1:fig-santai}
\end{bkfigure}

\label{ph:熱-1:joutai}状態変化では、分子の並び方が変わるだけで、分子の数は変わりません。だから、\textbf{質量は変わらず、体積が変わります}。水が氷になると、密度が約\qty{1000}{kg/m^3}から\qty{917}{kg/m^3}に小さくなるので、体積は $1000\div917\approx1.09$ 倍、約9\,\% 増えます。冬に水道管が破れるのは、このためです。

\phboxhead{とけている間、温度は変わらない}
氷を温めていくと、0\,℃でとけ始めます。とけている間は、熱を入れ続けても、温度は0\,℃のままです。全部とけて水になってから、ふたたび温度が上がり始めます。わくときも同じで、1気圧の水は、わいている間、約100\,℃のままです。\label{ph:熱-1:santai}小学校で学ぶ「水は約100\,℃でわき、0\,℃で氷になる。その間は温度が変わらない」は、このことです。

では、とけている間に入れた熱は、どこへ行ったのでしょう。氷の中の水の分子は、となりの分子と手をつなぐように引き合って、きちんと並んでいます（水素結合）。とけるには、そのつながりを切って、分子が動き回れるようにする必要があります。入れた熱は、\textbf{つながりを切る仕事に使われて、分子の動きを速くする（温度を上げる）分が残らない}のです。

\phboxhead{潜熱：$Q=mL$}
状態変化に使われ、温度の変化として表に出てこない熱を\textbf{潜熱}（せんねつ）といいます。\bkidx{せんねつ}{潜熱}「ひそんでいる熱」という意味です。\qty{1}{kg}あたりの潜熱を $L$〔J/kg〕とすると、質量 $m$ を状態変化させるのに要る熱は
\pheq{Q=mL}{J ＝ kg × J/kg}
\label{ph:熱-1:mL}です。固体をとかす潜熱を\textbf{融解熱}（$L_{\mathrm f}$）、液体を気体にする潜熱を\textbf{蒸発熱}（$L_{\mathrm v}$）といいます。\bkidx{ゆうかいねつ}{融解熱}\bkidx{じょうはつねつ}{蒸発熱}反対向きの変化（凍る・水蒸気が水にもどる）では、同じだけの熱が出てきます。

\begin{bkcenter}
\begin{bktable}{colspec={Q[l,m,wd=56mm] Q[c,m,wd=32mm] Q[l,m,wd=32mm]}, rowsep=1.2pt, row{1}={bg=bkTheme!10}}
潜熱 & $L$〔kJ/kg〕 & 出典 \\
氷の融解熱（0\,℃） & 334 & 理科年表 \\
水の蒸発熱（0.01\,℃） & 2501 & IF97 で計算 \\
水の蒸発熱（25\,℃） & 2442 & IF97 で計算 \\
水の蒸発熱（35\,℃） & 2418 & IF97 で計算 \\
水の蒸発熱（100\,℃） & 2257 & IF97 で計算 \\
\end{bktable}
\bktabcaption{水の潜熱。蒸発熱は、水の国際的な計算式 IAPWS-IF97 から計算した（巻末の蒸気表と同じ値）。}\label{ph:熱-1:tab-sennetsu}
\end{bkcenter}
蒸発熱は、融解熱の約7倍です。氷をとかすより、水を蒸発させるほうが、ずっと多くの熱が要ります。分子をばらばらに引きはなして、遠くまで飛べるようにするからです。また、水は100\,℃でなくても、ふだんの温度で少しずつ蒸発します。そのときにも、表の蒸発熱が要ります（低い温度ほど少し大きい）。

\begin{phyomi}[$Q=mL$]
「とかす（蒸発させる）熱は、とかす量に比例する」。温度の差が入っていないのは、状態変化の間は温度が変わらないから。$Q=mc\,\Delta T$ とは、使う場面が違う。
\end{phyomi}

\begin{phmijika}[汗と打ち水]
暑いとき、体は汗を出して冷やします。汗が皮ふの上で蒸発するとき、蒸発熱を体から奪うからです（例題7）。夏の道に水をまく打ち水も、水が蒸発するときに、地面と空気から熱を奪います。反対に、やかんの口から出る水蒸気でやけどしやすいのは、水蒸気が皮ふにふれて水にもどるとき、大きな蒸発熱を出すからです。
\end{phmijika}

\phboxhead{\bkhatenaref の答え}
氷水の中では、まわりから入った熱は、氷をとかす潜熱に使われ、温度を上げる分が残りません。それに、もし水の一部が0\,℃より温かくなっても、その水は、ふれている氷に熱を渡して、氷をとかしながら0\,℃にもどります。1気圧の下で、\textbf{氷と水がいっしょにいられる温度は0\,℃だけ}なのです。だから、氷が残っているかぎり、よくかき混ぜた氷水は0\,℃のままです。氷がとけきると、入った熱が水の温度を上げ始めます。

\begin{seikaku}
融点と沸点は、まわりの圧力で変わります。水の沸点は、気圧の低い富士山の頂上（約\qty{63}{kPa}）では約87\,℃に下がり（熱-8）、圧力なべの中（約\qty{200}{kPa}）では約120\,℃に上がります（熱-8・熱-10）。氷の融点は、圧力を上げるとわずかに下がりますが、約100気圧で約\qty{0.7}{K}ほどで、とても小さな変化です（熱-10）。また、食塩水のように何かがとけている水は、0\,℃より低い温度で凍ります。
\end{seikaku}

\begin{phmada}[氷の上は、なぜよく滑る？]
スケートの刃が氷の上をよく滑るのは、「刃の圧力で氷がとけて、水の膜ができるから」と説明されることがあります。でも、圧力でとける温度の下がり方は小さく（上の「せいかくに言うと」）、$-35$\,℃ のような寒さでもスケートやスキーができることを説明できません。今は、(1) 刃が擦れる摩擦の熱で、ごく薄い水の層ができること、(2) 氷の表面には、0\,℃より低くても、液体のように動く分子の薄い層がもともとあること、の2つが効いていると考えられています。でも、どちらがどれだけ効くのか、表面の層がどんなしくみでできるのかは、今も研究が続いています。身近な「滑る」の正体は、まだわかりきっていないのです。
\end{phmada}

\begin{column}{潜熱を見つけたジョゼフ・ブラック}
18世紀のなかばまで、熱と温度は、はっきり区別されていませんでした。スコットランドのジョゼフ・ブラック（1728〜1799年）は、1760年ごろ、氷を温めても、とけている間は温度が上がらないことに注目しました。そして、「温度計には表れないけれど、氷をとかすのに使われる熱がある」と考え、これを\textbf{潜熱}とよびました。ブラックは、同じ熱を入れても物によって温度の上がり方が違うこと（今の比熱の考え）も示しています。ただし、ブラックは熱を「熱素」という目に見えない物質だと考えていました。熱が物質ではなくエネルギーだとわかるのは、それから約80年後、ジュールの実験のころです（熱-2 ①）。
\end{column}

%% ---------------------------------------------------------------
\bksubsection{6}{熱の量を表にする：加熱曲線}
\label{ph:熱-1:sec6}
$-10$\,℃ の氷\qty{100}{g}を、1気圧のまま、100\,℃の水蒸気にするには、どれだけの熱が要るでしょう。温まる段と、状態が変わる段に分けて、表にします（比熱：氷 2.1、水 4.2\,kJ/(kg·K)。潜熱：融解 334、蒸発 2257\,kJ/kg）。
\begin{bkcenter}
\begin{bktable}{colspec={Q[l,m,wd=42mm] Q[l,m,wd=36mm] Q[r,m,wd=18mm] Q[r,m,wd=18mm]}, rowsep=1.2pt, row{1}={bg=bkTheme!10}, row{6}={font=\bfseries}}
段 & 式 & 熱〔kJ〕 & 合計〔kJ〕 \\
① 氷を $-10$→0\,℃ & $0.1\times2.1\times10$ & 2.1 & 2.1 \\
② 0\,℃で氷をとかす & $0.1\times334$ & 33.4 & 35.5 \\
③ 水を 0→100\,℃ & $0.1\times4.2\times100$ & 42.0 & 77.5 \\
④ 100\,℃で蒸発させる & $0.1\times2257$ & 225.7 & 303.2 \\
合計 & & 303.2 & \\
\end{bktable}
\bktabcaption{$-10$\,℃ の氷\qty{100}{g}を100\,℃の水蒸気にする熱。蒸発の段が全体の約74\,\% をしめる。}\label{ph:熱-1:tab-kanetsu}
\end{bkcenter}

\begin{bkfigure}
\begin{tikzpicture}
\begin{axis}[bkplot, width=118mm, height=52mm, xmin=0, xmax=335, ymin=-40, ymax=122, xtick={35.5,77.5,303.2}, xticklabels={35.5,77.5,303.2}, ytick={0,50,100}, xlabel={入れた熱（kJ）}, ylabel={$\theta$（℃）}, axis lines=left]
  \addplot[bkcurve] coordinates {(0,-10) (2.1,0) (35.5,0) (77.5,100) (303.2,100)};
  \addplot[bkcurve, dashed] coordinates {(303.2,100) (308.2,112)};
  \node[qxb lab, anchor=north west, text=bkBlue!80!black, fill=bkPaper] at (axis cs:9,-4) {② 0\,℃のまま、とける};
  \node[qxb lab, anchor=north west, text=bkBlue!80!black, fill=bkPaper] at (axis cs:9,-19) {① 氷が温まる（$-10$\,℃→0\,℃）};
  \node[qxb lab, anchor=south, text=bkBlue!80!black, fill=bkPaper] at (axis cs:165,103) {④ 100\,℃のまま、蒸発する};
  \node[qxb lab, anchor=west, text=bkBlue!80!black, fill=bkPaper] at (axis cs:62,48) {③ 水が温まる};
  \node[qxb lab, anchor=south east, text=bkGray, fill=bkPaper] at (axis cs:330,114) {水蒸気が温まる};
\end{axis}
\end{tikzpicture}
\bkfigcaption{加熱曲線。平らな所は状態変化（潜熱）、ななめの所は温まる所（比熱）。}\label{ph:熱-1:fig-kanetsu}
\end{bkfigure}

\phboxhead{グラフから読みとれること}
\begin{enumerate}[label={\bkcircled{\arabic*}}, leftmargin=2.4\zw, itemsep=1pt, topsep=2pt]
\item 平らな所の長さが、潜熱の大きさです。蒸発の段（④）は、とける段（②）の約7倍も長い。
\item ななめの所の傾きは $\Delta\theta/Q=1/(mc)$ です。氷の段（①）は、水の段（③）より約2倍急です。氷の比熱が水の約半分だからです。
\item 加える熱が1秒あたり一定（一定の電力）なら、横の軸は時間と同じです。やってみようのグラフで、氷が残っている間、温度が平らだったのは、②の段にいたからです。
\end{enumerate}

\phboxhead{表計算で、一定の電力で温める氷水を追う}
0\,℃の氷\qty{100}{g}と0\,℃の水\qty{200}{g}を、\qty{200}{W}のヒーターで温めます（熱は逃げないとする）。30秒ごとに入る熱は $200\times30=\qty{6000}{J}$ です。この熱を、まず氷をとかすのに使い、氷がなくなったら、残りの分で水の温度を上げます。質量は g、熱は J で計算します（融解熱 \qty{334}{J/g}、水の比熱 4.2\,J/(g·K)）。
\begin{bkcenter}
\bksheet{4}{24}{%
1 & 時刻（s） & 入った熱（J） & 残りの氷（g） & 温度（℃） \\
2 & 0 & 0 & 100.0 & 0.0 \\
3 & 30 & 6000 & 82.0 & 0.0 \\
4 & 60 & 12000 & 64.1 & 0.0 \\
5 & 90 & 18000 & 46.1 & 0.0 \\
6 & 120 & 24000 & 28.1 & 0.0 \\
7 & 150 & 30000 & 10.2 & 0.0 \\
8 & 180 & 36000 & 0.0 & 2.1 \\
9 & 210 & 42000 & 0.0 & 6.8 \\
10 & 240 & 48000 & 0.0 & 11.6 \\}
\bktabcaption{B3＝B2＋6000、C3＝MAX(0, C2－6000/334)、D3＝D2＋(6000－(C2－C3)*334)/(4.2*(300－C3))。}\label{ph:熱-1:tab-sheet}
\end{bkcenter}
C 列は「とけた分だけ氷が減る。ただし0より少なくはならない」、D 列は「6000\,J のうち、氷をとかすのに使わなかった分で、そのときの水（$300-\text{氷}$）の温度を上げる」という式です。150\,s から 180\,s の間に氷がなくなり、そこから先は、30秒ごとに $6000\div(4.2\times300)\approx\qty{4.8}{K}$ ずつ上がります。

\begin{phmodel}[一定の電力で温める氷水]
\phignore{入れ物・ヒーターが受けとる熱。まわりへ逃げる熱。}
\phassume{よくかき混ぜて、氷が残っている間は全体が0\,℃。ヒーターの熱は全部、中身に入る。}
\phlimit{かき混ぜないとき（場所で温度が違う）。温度が上がって、まわりへ逃げる熱が大きくなったとき（6-4 の形の式になる）。}
\end{phmodel}

%% ---------------------------------------------------------------
\bksubsection{7}{熱膨張：温めると、ふくらむ}
\label{ph:熱-1:sec7}
\phboxhead{長さが伸びる：線膨張（力-8 ⑦）}
力-8 ⑦で、長さ $L$ の棒の温度が $\Delta T$ 上がると、$\Delta L=\alpha L\,\Delta T$ だけ伸びることを学びました。$\alpha$ は線膨張係数で、鋼は約 $1.2\times10^{-5}$/K、アルミニウムは約 $2.3\times10^{-5}$/K、銅は約 $1.7\times10^{-5}$/K でした。ここでは、同じ3つの値を使って、\textbf{体積}の増え方を求めます。

\phboxhead{体積が増える：体膨張率 $\beta$}
\noindent\begin{minipage}[c]{0.36\linewidth}\centering
\begin{tikzpicture}[x=1mm, y=1mm]
  \filldraw[fill=bkBlue!20, draw=bkBlue!80!black, line width=0.6pt] (0,0) rectangle (20,20);
  \filldraw[fill=bkBlue!20, draw=bkBlue!80!black, line width=0.6pt] (0,20) -- (8,27) -- (28,27) -- (20,20) -- cycle;
  \filldraw[fill=bkBlue!30, draw=bkBlue!80!black, line width=0.6pt] (20,0) -- (28,7) -- (28,27) -- (20,20) -- cycle;
  \draw[bkRed!80!black, line width=0.7pt, dashed] (0,0) rectangle (23,23) (23,0) -- (31.6,7.5) -- (31.6,30.5) -- (23,23) (0,23) -- (8.6,30.5) -- (31.6,30.5);
  \draw[{Stealth[length=3pt]}-{Stealth[length=3pt]}, bkGray, line width=0.5pt] (0,-3) -- node[qxb lab, below] {$L$} (20,-3);
  \node[qxb lab, anchor=north, text=bkRed!75!black] at (24,-6.5) {伸びた辺 $L(1+\alpha\Delta T)$};
\end{tikzpicture}
\end{minipage}\hfill
\begin{minipage}[c]{0.6\linewidth}\raggedright
1辺 $L$ の立方体を考えます。温度が $\Delta T$ 上がると、どの辺も $L(1+\alpha\,\Delta T)$ になるので（向きによらず同じようにふくらむ物とする）、体積は
\[V'=L^3(1+\alpha\,\Delta T)^3=V(1+\alpha\,\Delta T)^3\]
です（$V=L^3$）。
\end{minipage}

\par\vspace{4pt}
$\alpha\,\Delta T$ は、たとえば鋼で\qty{100}{K}上げても $1.2\times10^{-3}$ と、1よりずっと小さい数です。そこで、$(1+h)^n\approx1+nh$（5-8④）を、$h=\alpha\,\Delta T$、$n=3$ で使うと、
\[V'\approx V(1+3\alpha\,\Delta T),\qquad \Delta V=V'-V\approx3\alpha V\,\Delta T\]
です。「1\,K あたり、体積の何倍だけ増えるか」を\textbf{体膨張率} $\beta$ といいます。\bkidx{たいぼうちょうりつ}{体膨張率}
\pheq{\Delta V=\beta V\,\Delta T,\qquad \beta\approx3\alpha}{m$^3$ ＝ 1/K × m$^3$ × K}
\label{ph:熱-1:3alpha}鋼なら $\beta\approx3.6\times10^{-5}$/K、アルミニウムなら $6.9\times10^{-5}$/K、銅なら $5.1\times10^{-5}$/K です。すてた項 $3(\alpha\,\Delta T)^2$ は、鋼を\qty{100}{K}上げたときで $3\times(1.2\times10^{-3})^2\approx4\times10^{-6}$ にすぎません。

\begin{phflow}[体膨張]
\phflowrow{$V'=V(1+\alpha\,\Delta T)^3$\\（どの向きも同じに伸びる）}{$\alpha\,\Delta T\ll1$}{$\Delta V\approx3\alpha V\,\Delta T$}
\end{phflow}

\phboxhead{液体と気体は、もっとふくらむ}
液体は、決まった形をもたないので、体膨張率 $\beta$ だけを使います。気体は、もっと大きくふくらみます（気体の体積と温度の関係は、熱-3 ②で式にします）。
\begin{bkcenter}
\begin{bktable}{colspec={Q[l,m,wd=44mm] Q[c,m,wd=34mm] Q[l,m,wd=38mm]}, rowsep=1.2pt, row{1}={bg=bkTheme!10}}
物（20\,℃前後） & $\beta$〔1/K〕 & 鋼とくらべると \\
空気（圧力一定） & $3.4\times10^{-3}$ & 約95倍 \\
ガソリン & $9.5\times10^{-4}$ & 約26倍 \\
水 & $2.1\times10^{-4}$ & 約6倍 \\
鋼（$3\alpha$） & $3.6\times10^{-5}$ & 1倍 \\
\end{bktable}
\bktabcaption{体膨張率。空気＞水＞金属の順に、桁で違う（空気の値は $1/T$ で、熱-3 ②）。}\label{ph:熱-1:tab-beta}
\end{bkcenter}
\label{ph:熱-1:bouchou}小学校で学ぶ「物は温めると体積が増え、冷やすと減る。増え方は空気＞水＞金属」は、この表のことです。液柱温度計は、液のふくらみを、細い管で長さに拡大して読む道具です。

\begin{phbaisuru}{温度の変化 $\Delta T$}
$\Delta V=\beta V\,\Delta T$ なので、増える体積も2倍です（$\beta$ がほぼ一定の範囲）。\qty{50}{L}のガソリンなら、\qty{10}{K}で約\qty{0.48}{L}、\qty{20}{K}で約\qty{0.95}{L}増えます（例題8）。
\end{phbaisuru}

\phboxhead{水は4\,℃でいちばん重い}
\noindent\begin{minipage}[c]{0.5\linewidth}\centering
\begin{tikzpicture}
\begin{axis}[bkplot, width=54mm, height=44mm, xmin=0, xmax=21, ymin=998.0, ymax=1000.25, xtick={0,4,8,12,16,20}, ytick={998.5,999,999.5,1000}, yticklabels={998.5,999.0,999.5,1000.0}, xlabel={$\theta$（℃）}, ylabel={密度（kg/m$^3$）}, axis lines=left]
  \addplot[bkcurve, smooth] coordinates {(0.1,999.85) (1,999.90) (2,999.94) (3,999.97) (4,999.97) (5,999.97) (6,999.94) (8,999.85) (10,999.70) (12,999.50) (14,999.25) (16,998.95) (18,998.60) (20,998.21)};
  \fill[bkRed] (axis cs:3.98,999.975) circle (1.8pt);
\end{axis}
\end{tikzpicture}
\end{minipage}\hfill
\begin{minipage}[c]{0.47\linewidth}\raggedright
ふつうの液体は、温度が下がると縮み、密度が大きくなります。ところが水は、約4\,℃（3.98\,℃）で密度がいちばん大きく、それより冷えると、かえってふくらみます（CoolProp 8.0 で計算）。0\,℃と4\,℃の間では、$\beta$ が負になるのです。
\end{minipage}

\begin{phmijika}[冬の池は、上から凍る]
冬に池の水が冷えると、冷えて重くなった水が底へしずみ、全体がかき混ぜられます。でも、4\,℃より冷えた水は軽いので、上にたまります。そのため、池は表面から凍り、底には4\,℃に近い水が残ります。氷は水より軽いので浮き、ふたのように下の水を冷えにくくします。魚が冬をこせるのは、このおかげです。
\end{phmijika}

\begin{seikaku}
液体や気体の体積は、温度だけでなく圧力でも変わります。体積を温度 $T$ と圧力 $p$ の関数 $V(T,\ p)$ とみると、全微分（8-3⑤）は
\[\dd V=\Bigl(\frac{\partial V}{\partial T}\Bigr)_p\dd T+\Bigl(\frac{\partial V}{\partial p}\Bigr)_T\dd p\]
です。体膨張率は、圧力を止めたときの1つめの項から
\[\beta=\frac{1}{V}\Bigl(\frac{\partial V}{\partial T}\Bigr)_p\]
と決めます。右下の $p$ は「圧力を止めて」という印です（8-3 の専門書）。
\end{seikaku}

%% ---------------------------------------------------------------
\bksubsection{8}{燃える量を数える：空気の量と発熱量}
\label{ph:熱-1:sec8}
\phboxhead{燃えるとは：原子の組みかえ}
ガスこんろやボイラーでは、燃料を燃やして熱を作ります。燃えるとは、燃料が空気の中の酸素と結びつく変化です。都市ガスのおもな成分のメタン（CH$_4$）なら、
\[\mathrm{CH_4}+2\,\mathrm{O_2}\ \longrightarrow\ \mathrm{CO_2}+2\,\mathrm{H_2O}\]
です。燃える前と後で、原子は組みかわるだけで、炭素 C・水素 H・酸素 O の原子の数はそれぞれ変わりません。これは、力-13 で学んだ「質量のバランス式」の、原子ごとのいちばん簡単な場合（生じる量も消える量も0）です。
\begin{bkcenter}
\begin{bktable}{colspec={Q[l,m,wd=24mm] Q[c,m,wd=34mm] Q[c,m,wd=34mm] Q[c,m,wd=16mm]}, rowsep=1.2pt, row{1}={bg=bkTheme!10}}
原子 & 前：CH$_4$＋2O$_2$ & 後：CO$_2$＋2H$_2$O & 同じ？ \\
C & 1 & 1 & ○ \\
H & 4 & $2\times2=4$ & ○ \\
O & $2\times2=4$ & $2+2\times1=4$ & ○ \\
\end{bktable}
\bktabcaption{メタンが燃えるときの、原子のバランス。}\label{ph:熱-1:tab-genshi}
\end{bkcenter}

\phboxhead{理論空気量と空気比}
気体は、同じ温度・同じ圧力なら、同じ体積に同じ数の分子をふくみます（アボガドロの法則。熱-3 ③）。だから、上の式の係数は、そのまま体積の比として使えます。メタン\qty{1}{m^3}を燃やすには、酸素が\qty{2}{m^3}要ります。燃焼の計算では、体積を0\,℃・1気圧（\qty{101.325}{kPa}）の体積に直して、単位を m$^3_{\mathrm N}$（ノルマル立方メートル）と書きます。\bkidx{のるまるりっぽうめーとる}{m$^3_{\mathrm N}$（ノルマル立方メートル）}

空気は、体積で約21\,\% が酸素、残りの約79\,\% がおもに窒素です。燃料をちょうど全部燃やすのに要る空気の量を\textbf{理論空気量}といいます。\bkidx{りろんくうきりょう}{理論空気量}メタン1\qxbmN{}では
\[A_0=\frac{\text{要る酸素}}{0.21}=\frac{2}{0.21}\approx9.5\qxbmN{}\]
です。本物のボイラーやこんろでは、燃え残りが出ないように、理論空気量より少し多めの空気を送ります。送った空気の量 $A$ と理論空気量 $A_0$ の比を\textbf{空気比} $\lambda$（ラムダ）といいます。\bkidx{くうきひ}{空気比}
\pheq{\lambda=\frac{A}{A_0}}{m$^3_{\mathrm N}$ ÷ m$^3_{\mathrm N}$：単位なし}
ガスのボイラーでは、$\lambda$ は 1.1〜1.3 くらいです。多すぎると、燃えるのに使われない空気まで温めて、煙突から捨ててしまうので、むだになります。

\phboxhead{高発熱量と低発熱量}
燃料\qty{1}{mol}（または\qty{1}{kg}、1\qxbmN{}）が燃えたときに出る熱を\textbf{発熱量}といいます。\bkidx{はつねつりょう}{発熱量}燃えてできた水（H$_2$O）を、液体の水にもどすところまで数えた値を\textbf{高発熱量}（総発熱量）、水蒸気のままにして、蒸発熱の分を数えない値を\textbf{低発熱量}（真発熱量）といいます。2つの差は、できた水の蒸発熱です。
\begin{bkcenter}
\begin{tikzpicture}[x=1mm, y=1mm]
  \fill[bkOrange!35] (0,0) rectangle (89.06,8);
  \fill[bkBlue!25] (80.26,0) rectangle (89.06,8);
  \draw[bkInk, line width=0.6pt] (0,0) rectangle (89.06,8);
  \draw[bkInk, line width=0.6pt] (80.26,0) -- (80.26,8);
  \draw[{Stealth[length=3pt]}-{Stealth[length=3pt]}, bkGray, line width=0.5pt] (0,11) -- node[qxb lab, above] {高発熱量 \qty{890.6}{kJ/mol}} (89.06,11);
  \draw[{Stealth[length=3pt]}-{Stealth[length=3pt]}, bkGray, line width=0.5pt] (0,-3) -- node[qxb lab, below] {低発熱量 \qty{802.6}{kJ/mol}} (80.26,-3);
  \node[qxb lab, anchor=west, align=left] at (91,4) {水 2\,mol の\\蒸発熱 88\,kJ};
\end{tikzpicture}
\end{bkcenter}
メタンでは、高発熱量が約\qty{890.6}{kJ/mol}、低発熱量が約\qty{802.6}{kJ/mol}（25\,℃）で、差の\qty{88}{kJ/mol}は、水 2\,mol の蒸発熱（25\,℃で 1\,mol あたり\qty{44.0}{kJ}）です（例題10）。日本の都市ガス（13A）の料金は、高発熱量で45\,MJ/m$^3_{\mathrm N}$を基準にしています。

\begin{phmijika}[潜熱を取りもどす給湯器]
ふつうのガス給湯器は、燃えてできた水蒸気を、そのまま煙突から出します。「潜熱回収型」の給湯器は、もう1つの熱交換器で排気を冷やし、水蒸気を水にもどすときに出る蒸発熱まで、お湯を温めるのに使います。そのぶん、同じお湯を作るのに要るガスが1割あまり少なくなります（そのかわりに、酸っぱい水（ドレン）が出るので、捨てる管が要ります）。
\end{phmijika}

\begin{phsaki}{燃焼工学}
この節では、燃える前と後の量だけを数えました。どれくらいの速さで燃えるか、炎の形はどう決まるか、すすや窒素酸化物（NO$_x$）がどうしてできるかは、化学反応の速さと、熱や物の流れ（熱-12）をいっしょに考える\textbf{燃焼工学}で学びます。専門の本では、何百もの化学反応を同時に計算する方法も使います。
\end{phsaki}

%% ---------------------------------------------------------------
\bksubsection{9}{熱をはかる：熱量計}
\label{ph:熱-1:sec9}
\noindent\begin{minipage}[c]{0.4\linewidth}\centering
\begin{tikzpicture}[x=1mm, y=1mm]
  \filldraw[fill=bkGray!25, draw=bkInk, line width=0.7pt] (0,0) rectangle (40,36);
  \fill[white] (4,4) rectangle (36,36);
  \draw[bkInk, line width=0.7pt] (4,36) -- (4,4) -- (36,4) -- (36,36);
  \fill[bkBlue!20] (8,6) rectangle (32,26);
  \draw[bkInk, line width=0.8pt] (8,30) -- (8,6) -- (32,6) -- (32,30);
  \filldraw[fill=bkGray!45, draw=bkInk, line width=0.7pt] (-2,36) rectangle (42,40);
  \filldraw[fill=bkOrange!40, draw=bkInk, line width=0.6pt] (16,10) rectangle (24,16);
  \draw[bkInk, line width=0.6pt] (20,16) -- (20,40);
  \qxbthermo{(27,12)}{36}{14}{bkRed}
  \draw[bkInk, line width=0.8pt] (11,8) -- (11,44) (8,8) -- (14,8);
  \node[qxb lab, anchor=east, fill=bkPaper] at (-1,22) {断熱材};
  \node[qxb lab, anchor=south] at (11,44.5) {かき混ぜ棒};
  \node[qxb lab, anchor=north] at (20,-1) {金属の試料と水};
\end{tikzpicture}
\end{minipage}\hfill
\begin{minipage}[c]{0.57\linewidth}\raggedright
物の比熱を測る道具を\textbf{熱量計}といいます。\bkidx{ねつりょうけい}{熱量計}断熱材で囲んだ入れ物に水を入れ、かき混ぜ棒と温度計を入れておきます。そこへ、温めた金属の試料を入れ、全体が同じ温度になったときの温度を読みます。外と熱をやりとりしないので、④の熱のバランス式が使えます。
\[m_{\text{金}}c_{\text{金}}(\theta_{\text{金}}-\theta_{\mathrm f})=(m_{\text{水}}c_{\text{水}}+C_{\text{器}})(\theta_{\mathrm f}-\theta_{\text{水}})\]
\end{minipage}

\par\vspace{4pt}
左辺は金属が出した熱、右辺は水と入れ物が受けとった熱です。$C_{\text{器}}$ は、入れ物・かき混ぜ棒・温度計を合わせた熱容量です。これを水の質量に直した値（$C_{\text{器}}\div c_{\text{水}}$）を\textbf{水当量}（すいとうりょう）といいます。\bkidx{すいとうりょう}{水当量}$C_{\text{器}}$ を忘れると、出した熱が水だけに入ったことになり、比熱を小さく見積もってしまいます（例題5）。

\begin{seikaku}
本物の熱量計でも、少しずつ熱は逃げます。そこで、試料を入れる前と後の温度の変わり方（6-4 のように、ゆっくり冷める）を時間で記録し、逃げた分を引いて直します。また、比熱の値そのものは、こうして\textbf{測って決める}もので、理論だけでは決まりません。気体や単純な固体の比熱の一部は、分子の動きから説明できるようになりました（熱-14・熱-15）が、液体の水のように分子どうしが複雑に引き合う物の比熱は、簡単な理論の式では求められません。工学では、測った値（巻末の物性値の表）を使います。
\end{seikaku}

\begin{taisetsu}
\def\tsl#1{\makebox[6.4\zw][l]{\sffamily\bfseries\small\color{\bkTT}#1}}
\tsl{温度}$T=\theta+273.15$（K と℃）。温度の差 $\Delta T$ は K でも℃でも同じ数\\[3pt]
\tsl{熱量}$Q=Pt$（電力一定）。熱容量 $C_{\mathrm{th}}=\dfrac{Q}{\Delta T}=mc$\\[3pt]
\tsl{比熱}$Q=mc\,\Delta T$（比熱が一定、状態変化なし）\\[3pt]
\tsl{バランス}外と熱をやりとりしないとき $\sum C_i\,\Delta T_i=0$。$T_{\mathrm f}=\dfrac{C_1T_1+C_2T_2}{C_1+C_2}$\\[3pt]
\tsl{潜熱}$Q=mL$。状態変化の間は温度が一定\\[3pt]
\tsl{熱膨張}$\Delta L=\alpha L\,\Delta T$、$\Delta V=\beta V\,\Delta T$、$\beta\approx3\alpha$\\[3pt]
\tsl{燃焼}理論空気量 $A_0=\text{要る酸素}\div0.21$、空気比 $\lambda=A/A_0$
\end{taisetsu}

\begin{kotoba}
\small
\bkword{熱平衡}{くっつけた物どうしの温度がそろい、熱が移らなくなった状態。}
\bkword{絶対温度}{絶対零度を0にした温度 $T$〔K〕。$T=\theta+273.15$。}
\bkword{熱容量}{物の温度を\qty{1}{K}上げるのに要る熱〔J/K〕。}
\bkword{比熱}{1\,kg あたりの熱容量〔J/(kg·K)〕。}
\bkword{潜熱}{状態変化に使われ、温度の変化に出てこない熱。融解熱・蒸発熱。}
\bkword{体膨張率}{\qty{1}{K}あたり、体積の何倍だけ増えるか〔1/K〕。}
\bkword{理論空気量・空気比}{燃料をちょうど燃やす空気の量と、送った空気との比。}
\bkword{高発熱量・低発熱量}{できた水の蒸発熱を数えるか、数えないか。}
\end{kotoba}

\bkkigocards{\bkkigocard{$\theta$}{シータ}{セルシウス温度〔℃〕。この巻では、式の中の $T$ はいつも絶対温度〔K〕で、℃の値は $\theta$ と書く。数学第6巻は温度を $T$〔℃〕と書いた}{Celsius temperature}{熱-1}}{\bkkigocard{$T$}{ティー}{絶対温度〔K〕。$T=\theta+273.15$。温度の差 $\Delta T$ は K でも℃でも同じ数}{thermodynamic temperature}{熱-1}}
\bkkigocards{\bkkigocard{$c$}{シー}{比熱〔J/(kg·K)〕。光の速さ $c_0$（熱-13）とは別}{specific heat capacity}{熱-1}}{\bkkigocard{$L$}{エル}{潜熱〔J/kg〕。融解熱 $L_{\mathrm f}$、蒸発熱 $L_{\mathrm v}$}{latent heat}{熱-1}}
\bkkigocards{\bkkigocard{$\beta$}{ベータ}{体膨張率〔1/K〕。線膨張係数は $\alpha$（力-8）}{volumetric expansion coefficient}{熱-1}}{\bkkigocard{$k_{\mathrm B}$}{ケー・ビー}{ボルツマン定数 \qty{1.380649e-23}{J/K}（決めた値）。いつも添字 B をつける（熱伝導率 $k$ などと区別）}{Boltzmann constant}{熱-1}}

%% ======================== 例題 ========================
\Needspace{10\baselineskip}\bksection{reidai}{例題でたしかめよう}
\noindent 水の比熱は 4.2\,kJ/(kg·K)、水\qty{1}{L}は\qty{1}{kg}とします。答えは2けたで書きます（例題1だけは3けたで計算します）。

\label{ph:熱-1:ex1}
\begin{phreidai}[氷水の氷は、全部とけるか（\bkhatenaref）]
断熱したコップに、0\,℃の氷\qty{100}{g}と、20\,℃の水\qty{400}{g}を入れて、よくかき混ぜました。氷は全部とけますか。最後の温度は何℃ですか。この例題は、出せる熱ととける熱がとても近いので、3けたで計算します（水の比熱は 4.18\,kJ/(kg·K)、融解熱は 334\,kJ/kg）。
\phbamen
温かい水が熱を出し、その熱で氷がとけます。水が出せる熱には限りがあります（氷があるかぎり、水は0\,℃より下がらない）。
\phmushi
コップと空気が受けとる熱を無視します（断熱）。氷ははじめから0\,℃とします。
\phzu
氷 $m_{\text{氷}}=\qty{0.100}{kg}$、水 $m_{\text{水}}=\qty{0.400}{kg}$、$\theta_{\text{水}}=20\,℃$、$L_{\mathrm f}=334$\,kJ/kg、$c=4.18$\,kJ/(kg·K)。
\phriyu
状態が変わるかどうかで式が変わるので、まず「氷を全部とかす熱」と「水が0\,℃まで下がって出せる熱」をくらべます（⑤の $Q=mL$ と③の $Q=mc\,\Delta T$）。そのあと、熱のバランス式（④）で最後の温度を決めます。
\phtateru
\[Q_{\text{とかす}}=m_{\text{氷}}L_{\mathrm f},\qquad Q_{\text{出せる}}=m_{\text{水}}c\,(20-0)\]
\phtoku
\[Q_{\text{とかす}}=0.100\times334=\qty{33.4}{kJ},\qquad Q_{\text{出せる}}=0.400\times4.18\times20=\qty{33.44}{kJ}\]
出せる熱のほうが\qty{0.04}{kJ}だけ多いので、氷は全部とけます。残った\qty{40}{J}が、とけた水もふくめた\qty{0.500}{kg}の水を温めます。
\[\Delta T=\frac{40}{0.500\times4180}\approx\qty{0.02}{K}\]
\phtashikame
大きさ：水が\qty{20}{K}下がって出す熱と、氷をとかす熱が、ほぼ同じです。氷\qty{1}{g}をとかす熱（\qty{334}{J}）で、水\qty{1}{g}を約\qty{80}{K}下げられるので、水が4倍で20\,℃なら、ちょうどつり合います（$4\times20=80$）。\\
極端な場合：氷を\qty{150}{g}にすると、とかす熱\qty{50.1}{kJ}が足りなくなり、氷が残って0\,℃になります。
\kotae 氷はほぼちょうど全部とけ、温度は約0\,℃（0.02\,℃）
\end{phreidai}
\begin{jibun}
断熱したコップに、0\,℃の氷\qty{50}{g}と30\,℃の水\qty{300}{g}を入れました。氷は全部とけますか。最後の温度を求めましょう（融解熱 334\,kJ/kg）。
\bkjans{とかす熱 $0.050\times334=\qty{16.7}{kJ}$ ＜ 水が0\,℃まで出せる熱 $0.300\times4.2\times30=\qty{37.8}{kJ}$ なので、全部とける。$0.300\times4.2\times(30-\theta)=16.7+0.050\times4.2\times\theta$ より $\theta\approx14\,℃$}
\end{jibun}

\label{ph:熱-1:ex2}
\begin{phreidai}[やかんでお湯をわかす熱]
やかんに水\qty{2.0}{L}を入れ、15\,℃から65\,℃まで温めます。水が受けとる熱を求めましょう。
\phbamen
状態は変わらず（水のまま）、温度だけが上がります。
\phmushi
やかんが受けとる熱、まわりへ逃げる熱、蒸発を無視します。
\phzu
$m=\qty{2.0}{kg}$、$c=4.2$\,kJ/(kg·K)、$\Delta T=65-15=\qty{50}{K}$。
\phriyu
状態変化がなく、比熱はほぼ一定（図\ref{ph:熱-1:fig-cw}）なので、$Q=mc\,\Delta T$（③）が使えます。温度の差は K でも℃でも同じ数（②）です。
\phtateru
\[Q=mc\,\Delta T\]
\phtoku
\[Q=2.0\times4.2\times50=\qty{420}{kJ}\]
\phtashikame
単位：kg × kJ/(kg·K) × K ＝ kJ。\\
大きさ：\qty{1.2}{kW}のこんろの熱が全部入るなら、$420\div1.2=\qty{350}{s}$（約6分）。台所での実感と合う桁です。\\
2倍にすると：水を\qty{4.0}{L}にすると、熱も2倍の\qty{840}{kJ}要ります。
\kotae 約\qty{420}{kJ}（\qty{4.2e5}{J}）
\end{phreidai}
\begin{jibun}
水\qty{1.5}{L}を20\,℃から100\,℃まで温める熱を求めましょう。
\bkjans{$1.5\times4.2\times80=\qty{504}{kJ}$ より、約\qty{5.0e2}{kJ}}
\end{jibun}

\label{ph:熱-1:ex3}
\begin{phreidai}[鉄なべの熱容量]
質量\qty{1.5}{kg}の鉄なべの熱容量を求めましょう。また、それは水何 g 分の熱容量ですか（鉄の比熱 0.45\,kJ/(kg·K)）。
\phbamen
なべも、中の食べ物といっしょに温まります。なべを温めるのに、どれだけの熱が要るかを知りたい場面です。
\phmushi
なべはどこも同じ鉄で、温度はどこも同じとします（取っ手などは考えない）。
\phzu
$m=\qty{1.5}{kg}$、$c_{\text{鉄}}=450\,J/(kg·K)$。比べる水の比熱 $c_{\text{水}}=4200\,J/(kg·K)$。
\phriyu
熱容量は「質量×比熱」（③）。同じ熱容量の水の質量は、熱容量を水の比熱でわれば出ます（⑨の水当量と同じ考え）。
\phtateru
\[C_{\mathrm{th}}=mc_{\text{鉄}},\qquad m_{\text{水}}=\frac{C_{\mathrm{th}}}{c_{\text{水}}}\]
\phtoku
\[C_{\mathrm{th}}=1.5\times450=\qty{675}{J/K},\qquad m_{\text{水}}=\frac{675}{4200}\approx\qty{0.16}{kg}\]
\phtashikame
単位：(J/K)÷(J/(kg·K)) ＝ kg。\\
大きさ：重い鉄なべ\qty{1.5}{kg}でも、熱容量は水\qty{160}{g}（コップ1ぱい弱）分しかありません。鉄の比熱が水の約 $\tfrac{1}{9}$ だからです。
\kotae 約\qty{680}{J/K}。水約\qty{160}{g}分
\end{phreidai}
\begin{jibun}
質量\qty{0.80}{kg}のアルミニウムのなべ（比熱 0.90\,kJ/(kg·K)）の熱容量と、それが水何 g 分かを求めましょう。
\bkjans{$0.80\times900=\qty{720}{J/K}$。$720\div4200\approx\qty{0.17}{kg}$ より、水約\qty{170}{g}分}
\end{jibun}

\label{ph:熱-1:ex4}
\begin{phreidai}[お湯と水を混ぜる]
80\,℃のお湯\qty{200}{g}に、20\,℃の水\qty{300}{g}を混ぜました。混ぜたあとの温度を求めましょう。
\phbamen
熱いお湯が出した熱を、冷たい水が受けとり、同じ温度になります。
\phmushi
入れ物とまわりの空気が受けとる熱、蒸発を無視します。
\phzu
お湯 $m_1=\qty{200}{g}$、$\theta_1=80\,℃$。水 $m_2=\qty{300}{g}$、$\theta_2=20\,℃$。最後の温度 $\theta_{\mathrm f}$。
\phriyu
外と熱をやりとりしないので、熱のバランス式（④）が使えます。同じ水どうしなので、比熱が約分できます。
\phtateru
\[m_1c(\theta_{\mathrm f}-80)+m_2c(\theta_{\mathrm f}-20)=0\]
\phtoku
$c$ でわって、$200(\theta_{\mathrm f}-80)+300(\theta_{\mathrm f}-20)=0$ です。
\[500\,\theta_{\mathrm f}=200\times80+300\times20=\num{22000}\qquad \theta_{\mathrm f}=44\,℃\]
\phtashikame
大きさ：20と80のまん中の50より、量の多い水の側（20）に寄っています（てんびん、図\ref{ph:熱-1:fig-tenbin}）。\\
熱で確かめる：お湯が出した熱 $0.200\times4.2\times36=\qty{30.24}{kJ}$、水が受けとった熱 $0.300\times4.2\times24=\qty{30.24}{kJ}$ で同じです。
\kotae 44\,℃
\end{phreidai}
\begin{jibun}
60\,℃のお湯\qty{300}{g}に、10\,℃の水\qty{200}{g}を混ぜたあとの温度を求めましょう。
\bkjans{$\dfrac{300\times60+200\times10}{500}=40\,℃$}
\end{jibun}

\label{ph:熱-1:ex5}
\begin{phreidai}[熱量計でアルミニウムの比熱を測る]
95\,℃に温めたアルミニウムの試料\qty{100}{g}を、熱量計の中の20\,℃の水\qty{200}{g}に入れたら、全体が26.8\,℃になりました。（1）入れ物の熱容量を無視して、アルミニウムの比熱を求めましょう。（2）入れ物・かき混ぜ棒・温度計の熱容量が\qty{60}{J/K}だとして、もう一度求めましょう。
\phbamen
試料が出した熱を、水（と入れ物）が受けとります。わからないのは試料の比熱です。
\phmushi
熱量計の外へ逃げる熱を無視します。試料を移す間に冷えることも無視します。
\phzu
試料 $m_{\text{金}}=\qty{0.100}{kg}$、$95\to26.8\,℃$（\qty{68.2}{K}下がる）。水 $m_{\text{水}}=\qty{0.200}{kg}$、$20\to26.8\,℃$（\qty{6.8}{K}上がる）。入れ物 $C_{\text{器}}=\qty{60}{J/K}$。
\phriyu
⑨の熱量計のバランス式です。「出した熱＝受けとった熱」を、わからない $c_{\text{金}}$ について解きます。
\phtateru
\[m_{\text{金}}c_{\text{金}}\times68.2=(m_{\text{水}}c_{\text{水}}+C_{\text{器}})\times6.8\]
\phtoku
（1）$C_{\text{器}}=0$：$c_{\text{金}}=\dfrac{0.200\times4200\times6.8}{0.100\times68.2}=\dfrac{5712}{6.82}\approx840\,J/(kg·K)$\\
（2）$C_{\text{器}}=\qty{60}{J/K}$：$c_{\text{金}}=\dfrac{(840+60)\times6.8}{6.82}=\dfrac{6120}{6.82}\approx900\,J/(kg·K)$
\phtashikame
表の値：アルミニウムの比熱は 0.90\,kJ/(kg·K)（表\ref{ph:熱-1:tab-hinetsu}）。（2）は表と合い、（1）は約7\,\% 小さく出ました。入れ物が受けとった熱（$60\times6.8=\qty{408}{J}$）を数え忘れると、試料が出した熱を小さく見積もってしまうからです。\\
向き：熱が逃げたり、入れ物の分を忘れたりすると、比熱はいつも\textbf{小さく}出ます。表より大きな値が出たら、温度の読み間違いを疑います。
\kotae （1）約840\,J/(kg·K)　（2）約900\,J/(kg·K)
\end{phreidai}
\begin{jibun}
90\,℃の銅の試料\qty{200}{g}を、熱量計の中の20\,℃の水\qty{150}{g}に入れたら、27.1\,℃になりました。入れ物などの熱容量を\qty{50}{J/K}として、銅の比熱を求めましょう。
\bkjans{$\dfrac{(0.150\times4200+50)\times7.1}{0.200\times62.9}=\dfrac{4828}{12.58}$ より、約380\,J/(kg·K)}
\end{jibun}

\label{ph:熱-1:ex6}
\begin{phreidai}[氷を水蒸気にする熱]
$-10$\,℃ の氷\qty{100}{g}を、1気圧で100\,℃の水蒸気にするのに要る熱を求めましょう。そのうち、蒸発に使われる熱は何\,\% ですか。
\phbamen
温まる段と、状態が変わる段が、交互に出てきます（⑥の加熱曲線）。
\phmushi
入れ物とまわりへの熱を無視します。比熱は、氷・水それぞれで一定とします。
\phzu
$m=\qty{0.100}{kg}$。比熱 $c_{\text{氷}}=2.1$、$c_{\text{水}}=4.2$\,kJ/(kg·K)。$L_{\mathrm f}=334$、$L_{\mathrm v}=2257$\,kJ/kg。
\phriyu
温まる段は $Q=mc\,\Delta T$（③）、状態が変わる段は $Q=mL$（⑤）です。段に分けてたします（表\ref{ph:熱-1:tab-kanetsu}）。
\phtateru
\[Q=mc_{\text{氷}}\times10+mL_{\mathrm f}+mc_{\text{水}}\times100+mL_{\mathrm v}\]
\phtoku
\[Q=2.1+33.4+42.0+225.7=\qty{303.2}{kJ},\qquad \frac{225.7}{303.2}\approx0.744\]
\phtashikame
大きさ：蒸発の熱は、氷から100\,℃の水にするまでの熱（\qty{77.5}{kJ}）の約3倍です。なべの水をわかし始めてから、全部蒸発させるまでには、わくまでの時間の何倍もかかるのは、このためです。\\
単位：どの段も kg × kJ/kg（または kJ/(kg·K) × K）＝ kJ。
\kotae 約\qty{3.0e2}{kJ}。蒸発が約74\,\%
\end{phreidai}
\begin{jibun}
$-20$\,℃ の氷\qty{200}{g}を、50\,℃の水にするのに要る熱を求めましょう。そのうち、とかすのに使われる熱は何\,\% ですか。
\bkjans{$8.4+66.8+42.0=\qty{117.2}{kJ}$ より、約\qty{1.2e2}{kJ}。とかす熱は $66.8\div117.2$ より約57\,\%}
\end{jibun}

\label{ph:熱-1:ex7}
\begin{phreidai}[汗で体はどれだけ冷える？]
運動をして、汗\qty{50}{g}が皮ふの上で蒸発しました。体から奪われる熱を求めましょう。それは、体重\qty{60}{kg}の人の体温を何 K 下げる分ですか。皮ふの温度での水の蒸発熱を\qty{2.42}{MJ/kg}（35\,℃、表\ref{ph:熱-1:tab-sennetsu}）、体の平均の比熱を 3.5\,kJ/(kg·K) とします。
\phbamen
汗（ほぼ水）が蒸発するとき、蒸発熱を体から奪います（⑤の身近な例）。
\phmushi
汗の熱は、全部体から奪われるとします（まわりの空気から奪う分を無視）。流れ落ちて蒸発しない汗は考えません。体の温度はどこも同じとします。
\phzu
汗 $m=\qty{0.050}{kg}$、$L_{\mathrm v}=\qty{2420}{kJ/kg}$。体 $M=\qty{60}{kg}$、$c=3.5$\,kJ/(kg·K)。
\phriyu
蒸発は状態変化なので $Q=mL$（⑤）。その熱で体の温度が下がる分は、$Q=Mc\,\Delta T$（③）を $\Delta T$ について解きます。
\phtateru
\[Q=mL_{\mathrm v},\qquad \Delta T=\frac{Q}{Mc}\]
\phtoku
\[Q=0.050\times2420=\qty{121}{kJ},\qquad \Delta T=\frac{121}{60\times3.5}\approx\qty{0.58}{K}\]
\phtashikame
大きさ：コップ $\tfrac{1}{4}$ ぱいほどの汗で、体温の約\qty{0.6}{K}分の熱を奪えます。運動で体が作る熱（数百 W）を、汗が運び出しているのです。\\
2倍にすると：汗の量が2倍なら、奪う熱も2倍です。ただし、湿度が高いと汗は蒸発しにくく（熱-9）、流れ落ちるだけで体は冷えません。
\kotae 約\qty{1.2e2}{kJ}。体温の約\qty{0.58}{K}分
\end{phreidai}
\begin{jibun}
体重\qty{50}{kg}の人の汗\qty{100}{g}が蒸発したとき、奪われる熱と、それが体温の何 K 分かを求めましょう（蒸発熱 2.42\,MJ/kg、体の比熱 3.5\,kJ/(kg·K)）。
\bkjans{$0.100\times2420=\qty{242}{kJ}$。$242\div(50\times3.5)$ より、約\qty{1.4}{K}分}
\end{jibun}

\label{ph:熱-1:ex8}
\begin{phreidai}[ガソリンは、温まるとどれだけ増える？]
朝、10\,℃のガソリン\qty{50}{L}を、車の燃料タンクいっぱいに入れました。昼に、ガソリンの温度が30\,℃まで上がりました。ガソリンの体積は何 L 増えますか。ガソリンの体膨張率を $9.5\times10^{-4}$/K とします。また、鋼のタンク（線膨張係数 $1.2\times10^{-5}$/K）も同じだけ温まるとき、タンクの容積はどれだけ増えますか。
\phbamen
液体は温まるとふくらみます（⑦）。タンクもふくらみますが、ふくらみ方が違います。
\phmushi
ガソリンの蒸発、タンクの中の空気の部分を無視します。$\beta$ は10\,℃〜30\,℃で一定とします。
\phzu
$V=\qty{50}{L}$、$\Delta T=\qty{20}{K}$、$\beta_{\text{ガ}}=9.5\times10^{-4}$/K、タンク $\beta_{\text{鋼}}\approx3\alpha=3.6\times10^{-5}$/K。
\phriyu
体積の増え方は $\Delta V=\beta V\,\Delta T$（⑦）。鋼のタンクの容積は、中の空間の形のまま、鋼と同じ割合でふくらむので、$\beta\approx3\alpha$ を使います。
\phtateru
\[\Delta V_{\text{ガ}}=\beta_{\text{ガ}}V\,\Delta T,\qquad \Delta V_{\text{タ}}=3\alpha V\,\Delta T\]
\phtoku
\[\Delta V_{\text{ガ}}=9.5\times10^{-4}\times50\times20=\qty{0.95}{L}\]
\[\Delta V_{\text{タ}}=3.6\times10^{-5}\times50\times20=\qty{0.036}{L}\]
\phtashikame
くらべる：タンクの増え方はガソリンの約 $\tfrac{1}{26}$ で、差の約\qty{0.91}{L}があふれ出そうとします。燃料タンクやガソリンの携行缶に「満タンに入れすぎない」と書いてあるのは、このためです。\\
大きさ：$\beta\,\Delta T=0.019$ で、体積の約2\,\% です。
\kotae ガソリンは約\qty{0.95}{L}増える。タンクの容積は約\qty{0.036}{L}しか増えない
\end{phreidai}
\begin{jibun}
\qty{40}{L}のガソリンの温度が\qty{15}{K}上がると、体積は何 L 増えますか（$\beta=9.5\times10^{-4}$/K）。
\bkjans{$9.5\times10^{-4}\times40\times15=\qty{0.57}{L}$}
\end{jibun}

\label{ph:熱-1:ex9}
\begin{phreidai}[メタンを燃やす空気の量]
メタン 1\qxbmN を燃やすのに要る理論空気量を求めましょう。また、空気比1.2 で燃やしたとき、送る空気の量と、燃えたあとのガスから水を除いたもの（乾いた排ガス）にふくまれる酸素の割合（体積で）を求めましょう。
\phbamen
ボイラーで、少し多めの空気を送ってメタンを燃やします。余った酸素は、排ガスに残ります。
\phmushi
燃え残り（一酸化炭素など）はなく、全部が CO$_2$ と H$_2$O になるとします。空気は酸素21\,\%・窒素79\,\%（体積）とし、窒素は燃えないとします。
\phzu
CH$_4+2$O$_2\to$ CO$_2+2$H$_2$O。理論空気量 $A_0$、空気比 $\lambda=1.2$、送る空気 $A=\lambda A_0$。
\phriyu
原子のバランス（⑧）から、要る酸素は 2\qxbmN{}。気体の体積の比は分子の数の比なので、そのまま体積で計算できます。排ガスは、原子ごとに「入った分＝出た分」で数えます。
\phtateru
\[A_0=\frac{2}{0.21},\qquad A=1.2\,A_0\]
乾いた排ガス：CO$_2$ 1\qxbmN{}、窒素 $0.79A$、余った酸素 $0.21A-2$。
\phtoku
\[A_0=\frac{2}{0.21}\approx9.52\qxbmN{},\qquad A=1.2\times9.52\approx11.4\qxbmN{}\]
余った酸素 $0.21\times11.43-2=0.40$\qxbmN{}、窒素 $0.79\times11.43=9.03$\qxbmN{}。
\[\text{酸素の割合}=\frac{0.40}{1+9.03+0.40}=\frac{0.40}{10.43}\approx0.038\]
\phtashikame
大きさ：メタン1に対して、空気は約10倍要ります。空気の大部分は燃えない窒素で、その分の体積も温めてしまいます（空気比を上げすぎない理由）。\\
使い道：本物のボイラーでは、排ガスの酸素の割合を測って、空気比が適切かを確かめています。
\kotae 理論空気量 約9.5\qxbmN{}、送る空気 約11\qxbmN{}、酸素 約3.8\,\%
\end{phreidai}
\begin{jibun}
プロパン（C$_3$H$_8+5$O$_2\to3$CO$_2+4$H$_2$O）1\qxbmN の理論空気量と、空気比1.2 のときの空気の量を求めましょう。
\bkjans{$A_0=5\div0.21\approx23.8$\qxbmN より約24\qxbmN{}。$1.2\times23.8$ より約29\qxbmN{}}
\end{jibun}

\label{ph:熱-1:ex10}
\begin{phreidai}[高発熱量と低発熱量の差]
メタンの高発熱量は\qty{890.6}{kJ/mol}、低発熱量は\qty{802.6}{kJ/mol}です。差を求め、それがメタン 1\,mol が燃えてできる水の蒸発熱と同じことを確かめましょう（水の蒸発熱は25\,℃で\qty{44.0}{kJ/mol}）。また、高発熱量を1\,kg あたりに直しましょう（メタンのモル質量 \qty{16.04}{g/mol}）。
\phbamen
燃えてできた水を、液体にもどすまで数えるか（高）、水蒸気のままにするか（低）の違いです。
\phmushi
燃える前と後の温度は、どちらも25\,℃とします。
\phzu
CH$_4+2$O$_2\to$ CO$_2+2$H$_2$O で、水は 2\,mol できる。
\phriyu
2つの発熱量の差は、できた水の潜熱（⑤の $Q=mL$ を mol で書いたもの）です。1\,kg あたりに直すのは、単位の換算（2-10）です。
\phtateru
\[\Delta=890.6-802.6,\qquad 2\times44.0,\qquad \frac{890.6\ \text{kJ/mol}}{0.01604\ \text{kg/mol}}\]
\phtoku
\[\Delta=\qty{88.0}{kJ/mol},\qquad 2\times44.0=\qty{88.0}{kJ/mol}\]
\[\frac{890.6}{0.01604}\approx\num{55500}\ \text{kJ/kg}=\qty{55.5}{MJ/kg}\]
\phtashikame
同じ値：差の\qty{88.0}{kJ}は、水 2\,mol の蒸発熱とちょうど同じです。高発熱量の約1割が、水蒸気の潜熱です。\\
単位：(kJ/mol)÷(kg/mol) ＝ kJ/kg。
\kotae 差は\qty{88}{kJ/mol}（水 2\,mol の蒸発熱と同じ）。高発熱量は約\qty{56}{MJ/kg}（\qty{55.5}{MJ/kg}）
\end{phreidai}
\begin{jibun}
水素（H$_2+\tfrac{1}{2}$O$_2\to$ H$_2$O）の高発熱量は\qty{285.8}{kJ/mol}です。低発熱量を求め、両方を1\,kg あたりに直しましょう（水素のモル質量 \qty{2.016}{g/mol}、水の蒸発熱\qty{44.0}{kJ/mol}）。
\bkjans{低発熱量 $285.8-44.0=\qty{241.8}{kJ/mol}$。1\,kg あたり、高 約\qty{142}{MJ/kg}、低 約\qty{120}{MJ/kg}}
\end{jibun}

\label{ph:熱-1:ex11}
\begin{phreidai}[ガス給湯器の能力]
ガス給湯器で、15\,℃の水道の水を、毎分\qty{12}{L}ずつ40\,℃のお湯にします。お湯が受けとる熱は1秒あたり何 kW ですか。給湯器の効率（ガスの熱のうち、お湯に入る割合）を0.95 として、1秒あたりに要るガスの熱と、1時間に使うガスの量（13A：45\,MJ/m$^3_{\mathrm N}$）を求めましょう。
\phbamen
水が流れながら温まります。1秒あたりの量で数えます（記号の上の点は「1秒あたり」の印で、力-13 の $\dot m$ と同じです。$\dot V$ は体積流量、$\dot Q$ は1秒あたりの熱）。
\phmushi
配管から逃げる熱を、効率の中に入れます。水の密度は\qty{1}{kg/L}とします。
\phzu
$\dot V=\qty{12}{L/min}$、$\dot m=\qty{12}{kg/min}=\qty{0.20}{kg/s}$、$\Delta T=40-15=\qty{25}{K}$、$\eta=0.95$。
\phriyu
1秒に通る水 $\dot m$ が、それぞれ $c\,\Delta T$ ずつ熱を受けとるので、$Q=mc\,\Delta T$（③）を1秒あたりにした $\dot Q=\dot mc\,\Delta T$ を使います。ガスの熱は、効率でわります（力-4 ⑧）。
\phtateru
\[\dot Q_{\text{湯}}=\dot m\,c\,\Delta T,\qquad \dot Q_{\text{ガス}}=\frac{\dot Q_{\text{湯}}}{\eta}\]
\phtoku
\[\dot Q_{\text{湯}}=0.20\times4.2\times25=\qty{21}{kW},\qquad \dot Q_{\text{ガス}}=\frac{21}{0.95}\approx\qty{22}{kW}\]
1時間のガスの熱は $22.1\times3600\approx\num{79600}$\,kJ $=\qty{79.6}{MJ}$ なので、$79.6\div45\approx1.8$\qxbmN です。
\phtashikame
大きさ：\qty{21}{kW}は、電気ケトル（約\qty{1.2}{kW}）の約18台分です。お湯を「流しながら」作るには、大きな熱が要ります（工学ではこう使う）。\\
単位：kg/s × kJ/(kg·K) × K ＝ kJ/s ＝ kW。
\kotae お湯に約\qty{21}{kW}、ガスの熱 約\qty{22}{kW}、ガス 約1.8\qxbmN{}/h
\end{phreidai}
\begin{jibun}
シャワーで、15\,℃の水を毎分\qty{10}{L}ずつ40\,℃にします。お湯が受けとる熱（kW）と、効率0.80 の給湯器で要るガスの熱を求めましょう。
\bkjans{$\tfrac{10}{60}\times4.2\times25=\qty{17.5}{kW}$ より約\qty{18}{kW}。ガスの熱 $17.5\div0.80$ より約\qty{22}{kW}}
\end{jibun}

\label{ph:熱-1:ex12}
\begin{phreidai}[電気ポットの保温に要る電力]
電気ポットに90\,℃のお湯\qty{2.0}{L}が入っています。部屋は20\,℃です。ヒーターを切ると、10分で温度が\qty{2.0}{K}下がりました。（1）お湯から1秒あたりに逃げる熱、（2）90\,℃に保つのに要るヒーターの電力、（3）ポットの熱抵抗 $R_{\mathrm{th}}$ と時定数 $\tau$ を求めましょう。
\phbamen
保温とは、逃げる熱と同じだけの熱を、ヒーターで入れ続けることです（6-4⑥の「入る熱＝逃げる熱」）。
\phmushi
10分の間は、お湯と部屋の温度差（約\qty{70}{K}）がほとんど変わらないので、逃げる熱は一定とします。ポットの入れ物の熱容量と、湯気で出る熱は、逃げる熱にふくめて考えます。
\phzu
$C_{\mathrm{th}}=mc=2.0\times4200=\qty{8400}{J/K}$、10分＝\qty{600}{s}で $\Delta T=-\qty{2.0}{K}$、温度差 $90-20=\qty{70}{K}$。
\phriyu
ヒーターを切ったときは、熱のバランス式（④）が「たまる熱＝－逃げる熱」になります。保温のときは、たまる熱が0なので「入る熱＝逃げる熱」。熱抵抗と時定数は 6-4⑥ の決め方です。
\phtateru
\[\dot Q_{\text{逃}}=\frac{C_{\mathrm{th}}\,|\Delta T|}{t},\qquad P=\dot Q_{\text{逃}},\qquad R_{\mathrm{th}}=\frac{\text{温度差}}{\dot Q_{\text{逃}}},\qquad \tau=C_{\mathrm{th}}R_{\mathrm{th}}\]
\phtoku
\[\dot Q_{\text{逃}}=\frac{8400\times2.0}{600}=\qty{28}{W},\qquad P=\qty{28}{W}\]
\[R_{\mathrm{th}}=\frac{70}{28}=\qty{2.5}{K/W},\qquad \tau=8400\times2.5=\qty{21000}{s}\approx\qty{5.8}{h}\]
\phtashikame
大きさ：保温の電力\qty{28}{W}は、お湯をわかす電力（約\qty{1}{kW}）の約 $\tfrac{1}{35}$ です。それでも1日では $28\times24\approx\qty{670}{Wh}$ になります。\\
6-4 とくらべる：6-4⑥ のカップのお湯（\qty{1050}{J/K}、$R_{\mathrm{th}}\approx\qty{1.1}{K/W}$、$\tau=20$ 分）より、ポットは熱容量が大きく、熱抵抗も大きい（逃げにくい）ので、時定数が十数倍長くなります。
\kotae （1）約\qty{28}{W}　（2）約\qty{28}{W}　（3）$R_{\mathrm{th}}=\qty{2.5}{K/W}$、$\tau\approx\qty{2.1e4}{s}$（約5.8時間）
\end{phreidai}
\begin{jibun}
80\,℃のお湯\qty{1.0}{L}が入った電気ポットのヒーターを切ると、10分で\qty{3.0}{K}下がりました（部屋は20\,℃）。逃げる熱と、80\,℃に保つのに要る電力を求めましょう。
\bkjans{$4200\times3.0\div600=\qty{21}{W}$。保つ電力も約\qty{21}{W}}
\end{jibun}

%% ======================== まちがえやすいところ ========================
\Needspace{11\baselineskip}\bksection{tsumazuki}{まちがえやすいところをチェック}
\begin{tsumazuki}{20\,℃の2倍は40\,℃？}
\bkng{20\,℃から40\,℃になったので、2倍熱くなった。}
\bkok{絶対温度では\qty{293}{K}から\qty{313}{K}で、約1.07倍。温度の比は K で考える。}
℃の0は「水が凍る温度」で、エネルギーが0の点ではありません。比や $T$ そのものが効く式には、K を入れます（②）。
\end{tsumazuki}

\begin{tsumazuki}{温度の差も273をたす？}
\bkng{20\,℃から35\,℃への変化は、$35+273=308$ で、\qty{308}{K}上がった。}
\bkok{温度の差は K でも℃でも同じ数。\qty{15}{K}上がった。}
273をたすのは、温度そのものを K に直すときだけです。差をとると273は消えます（②）。
\end{tsumazuki}

\begin{tsumazuki}{0\,℃の氷をとかす熱に $mc\,\Delta T$？}
\bkng{0\,℃の氷が0\,℃の水になるので、$\Delta T=0$ で、熱は要らない。}
\bkok{状態変化には潜熱 $Q=mL$ が要る。氷\qty{1}{kg}なら\qty{334}{kJ}。}
$Q=mc\,\Delta T$ は「状態が変わらないで温度だけ変わる」ときの式です（③⑤）。加熱曲線の平らな所では $Q=mL$ を使います。
\end{tsumazuki}

\begin{tsumazuki}{比熱の kJ と、質量の g}
\bkng{水\qty{500}{g}を\qty{10}{K}上げる熱は $500\times4.2\times10=\num{21000}$\,kJ。}
\bkok{4.2\,kJ/(kg·K) なら、質量は kg で入れて $0.500\times4.2\times10=\qty{21}{kJ}$。}
4.2\,kJ/(kg·K)＝4.2\,J/(g·K) です。kJ と kg、J と g の組を、そろえて使います（2-10）。
\end{tsumazuki}

\begin{tsumazuki}{混ぜた温度は、2つの温度のまん中？}
\bkng{80\,℃のお湯と20\,℃の水を混ぜたら、いつも50\,℃。}
\bkok{量（熱容量）で重みをつけた平均。量が同じときだけ50\,℃。}
最後の温度は、熱容量の大きいほうに寄ります（④のてんびん）。
\end{tsumazuki}

%% ======================== 工学ではこう使う ========================
\Needspace{18\baselineskip}\bksection{phgear}{工学ではこう使う}
\begin{phkougaku}[給湯器の「号数」]
ガス給湯器のカタログには「20号」「24号」のような\textbf{号数}が書いてあります。これは、「水の温度を\qty{25}{K}上げたお湯を、1分間に何 L 出せるか」を表す数です。1号は、毎分\qty{1}{L}の水を\qty{25}{K}上げる能力なので、
\[\frac{1}{60}\ \text{kg/s}\times4.2\ \text{kJ/(kg·K)}\times\qty{25}{K}\approx\qty{1.75}{kW}\]
です（比熱を 4.186 として、約\qty{1.74}{kW}とする説明もあります）。例題11の給湯器（毎分\qty{12}{L}を\qty{25}{K}上げる）は、ちょうど12号です。冬に水道の水が5\,℃まで冷えると、40\,℃のお湯を作るには\qty{35}{K}上げる必要があり、同じ給湯器でも、出せるお湯は $12\times\tfrac{25}{35}\approx\qty{8.6}{L/min}$ に減ります。家族が多い家では、冬のシャワーと台所を同時に使えるよう、24号などの大きい給湯器を選びます。

\begin{phketa}[熱の量の目安]
\phketaline{水\qty{1}{L}を\qty{1}{K}上げる熱}{約\qty{4.2}{kJ}}
\phketaline{氷\qty{1}{kg}をとかす熱}{\qty{334}{kJ}}
\phketaline{水\qty{1}{kg}を蒸発させる熱}{約\qty{2.3}{MJ}〜\qty{2.5}{MJ}}
\phketaline{電気ケトル・24号のガス給湯器}{約\qty{1.2}{kW}・約\qty{42}{kW}}
\phketaline{人が出す熱（1日2000\,kcal。熱-2 例題2）}{約\qty{100}{W}}
\phketaline{都市ガス13A 1\qxbmN の熱（高発熱量）}{\qty{45}{MJ}}
\phboxhead{答えが桁で合っているか}
お湯をわかすのは数百 kJ、おふろ1ぱい（約\qty{200}{L}を\qty{25}{K}）で約\qty{21}{MJ}です。ふだんの家庭のお湯の計算で MJ をこえる答えが出たら、量が大きい（おふろ・給湯）のか、g と kg を取り違えたのかを確かめます。
\end{phketa}
\end{phkougaku}
\phstation{伝熱工学,燃焼工学,計測工学,空気調和}{熱容量、比熱、潜熱、水当量、理論空気量、空気比、高発熱量・低発熱量、号数}



%% ======================== 章末問題 ========================
\bkendproblems
\bkproblemgroup{基本}{1}
\bkproblem 温度を換算しましょう。（1）25\,℃は何 K ですか。（2）液体窒素の沸点\qty{77}{K}は何℃ですか。（3）18\,℃から33\,℃への温度の変化は何 K ですか。（4）68\,$^{\circ}$F は何℃ですか。
\bkpans{（1）\nobreak 約\qty{298}{K}（298.15）。（2）\nobreak 約 $-196$\,℃。（3）\nobreak \qty{15}{K}。（4）\nobreak $(68-32)\div1.8=20\,℃$}
\bkproblem 水\qty{3.0}{L}を10\,℃から60\,℃まで温める熱を求めましょう。また、\qty{1.5}{kW}のヒーターの熱が全部入るとすると、何分かかりますか。
\bkpans{$3.0\times4.2\times50=\qty{630}{kJ}$。$630\div1.5=\qty{420}{s}$ より、\qty{7.0}{min}}
\bkproblem 水\qty{0.80}{L}を、\qty{1000}{W}の電気ケトルで20\,℃から100\,℃にします。熱が逃げないとして、かかる時間を求めましょう（6-3③と同じ計算）。
\bkpans{$C_{\mathrm{th}}=\qty{3360}{J/K}$、熱 $3360\times80=\qty{268800}{J}$。$268800\div1000\approx\qty{269}{s}$ より、約\qty{4.5}{min}}
\bkproblem 70\,℃のお湯\qty{300}{g}と10\,℃の水\qty{100}{g}を混ぜたあとの温度を求めましょう。
\bkpans{$\dfrac{300\times70+100\times10}{400}=55\,℃$}
\bkproblem （1）0\,℃の氷\qty{250}{g}をとかす熱、（2）100\,℃の水\qty{250}{g}を全部蒸発させる熱を求めましょう（融解熱 334、蒸発熱 2257\,kJ/kg）。
\bkpans{（1）\nobreak 約\qty{84}{kJ}（83.5）。（2）\nobreak 約\qty{5.6e2}{kJ}（564）}
\bkproblem 長さ\qty{2.0}{m}の銅の棒の温度が\qty{50}{K}上がりました（$\alpha=1.7\times10^{-5}$/K）。（1）伸びを mm で求めましょう。（2）同じ銅の体積 \qty{1000}{cm^3} のかたまりでは、体積は何 cm$^3$ 増えますか（$\beta\approx3\alpha$）。
\bkpans{（1）\nobreak $1.7\times10^{-5}\times2.0\times50=\qty{1.7e-3}{m}$ より\qty{1.7}{mm}。（2）\nobreak $5.1\times10^{-5}\times1000\times50$ より約\qty{2.6}{cm^3}}
\bkproblem メタン 2.0\qxbmN を、空気比1.1 で燃やします。送る空気の量を求めましょう。
\bkpans{$2.0\times\dfrac{2}{0.21}\times1.1$ より、約21\qxbmN{}}
\bkproblemgroup{チャレンジ}{4}
\bkproblem 断熱した入れ物の中の、30\,℃の水\qty{400}{g}に、$-10$\,℃ の氷\qty{200}{g}を入れました。（1）氷は全部とけますか。（2）最後の温度と、残る氷の質量を求めましょう（氷の比熱 2.1、水の比熱 4.2\,kJ/(kg·K)、融解熱 334\,kJ/kg）。
\bkpans{（1）\nobreak 氷を0\,℃にしてとかす熱 $4.2+66.8=\qty{71.0}{kJ}$ ＞ 水が0\,℃まで出せる熱 $\qty{50.4}{kJ}$ なので、全部はとけない。（2）\nobreak 0\,℃。とける氷 $(50.4-4.2)\div334\approx\qty{0.138}{kg}$ なので、残る氷は約\qty{62}{g}}
\bkproblem 0\,℃の氷\qty{150}{g}と0\,℃の水\qty{350}{g}を、\qty{300}{W}のヒーターで温めます（熱は逃げない、よくかき混ぜる）。（1）氷がとけきるまでの時間、（2）そこから50\,℃になるまでの時間を求め、（3）時間と温度のグラフの形を説明しましょう。
\bkpans{（1）\nobreak $0.150\times334000\div300=\qty{167}{s}$。（2）\nobreak $0.500\times4200\times50\div300=\qty{350}{s}$。（3）\nobreak はじめの約\qty{167}{s}は0\,℃で平ら、そのあと1秒に約\qty{0.14}{K}ずつ直線で上がり、合計約\qty{517}{s}（約8.6分）で50\,℃}
\bkproblem おふろに、15\,℃の水道の水\qty{200}{L}を42\,℃にしてためます。（1）お湯が受けとる熱、（2）効率0.95 の給湯器で使う都市ガス（45\,MJ/m$^3_{\mathrm N}$）の量を求めましょう。
\bkpans{（1）\nobreak $200\times4.2\times27=\num{22680}$\,kJ より、約\qty{23}{MJ}。（2）\nobreak $22.68\div0.95\div45$ より、約0.53\qxbmN{}}
\bkproblem アルミニウムの立方体（$\alpha=2.3\times10^{-5}$/K）の温度を\qty{200}{K}上げます。体積の増える割合を、（1）$(1+\alpha\,\Delta T)^3-1$ で正確に、（2）$3\alpha\,\Delta T$ で近似して求め、（3）近似の誤差が何\,\% かを答えましょう。
\bkpans{（1）\nobreak $1.0046^3-1\approx0.013864$。（2）\nobreak $0.0138$。（3）\nobreak 近似は正確な値より約0.46\,\% 小さい（すてた項は $3(\alpha\Delta T)^2+(\alpha\Delta T)^3\approx6.4\times10^{-5}$）}

\begin{furikaeri}
\begin{checklist}
\item 氷がとけている間、温度が変わらないわけを、潜熱で説明できる。\bkfurref{ph:熱-1:sec5}
\item $Q=mc\,\Delta T$ と熱のバランス式で、温まり方と混ぜたときの温度を求められる。\bkfurref{ph:熱-1:sec4}
\item 加熱曲線を読み、熱膨張の大きさを見積もれる。\bkfurref{ph:熱-1:sec6}
\item 燃料が燃えるときの空気の量と発熱量を、原子のバランスで求められる。\bkfurref{ph:熱-1:sec8}
\end{checklist}
\end{furikaeri}
\bkfukidashi{mina}{温度は「熱い・冷たいの度合い」、熱は「移っていくエネルギーの量」なんだね。氷水が0\,℃のままなのは、入った熱が氷をとかすのに使われていたからだ。}
